%%%%%%%%%%%%%%%%%%%%%%%%%%%%%%%%%%%%%%%%%%%%%%%%%%%%%%%%%%%%
% 11.3 NUMERICAL OPTIMIZATION
% The Composition of Advanced Mathematics
%%%%%%%%%%%%%%%%%%%%%%%%%%%%%%%%%%%%%%%%%%%%%%%%%%%%%%%%%%%%

\section{Numerical Optimization}
\label{sec:numerical-optimization}

\prerequisite{
Optimization fundamentals, nonlinear programming, convex optimization,
multivariable calculus, gradients and Hessians, numerical analysis,
numerical linear algebra, and basic algorithmic complexity.
}

\learningobjectives{
By the end of this section, the reader should be able to:
\begin{itemize}
    \item explain the goals of numerical optimization;
    \item distinguish analytical and numerical optimization;
    \item define iterates, search directions, and step lengths;
    \item understand descent directions;
    \item understand line-search methods;
    \item distinguish exact and inexact line searches;
    \item apply Armijo sufficient decrease;
    \item understand Wolfe and strong Wolfe conditions;
    \item understand gradient descent;
    \item analyze steepest descent on quadratic functions;
    \item understand the role of conditioning in convergence;
    \item derive Newton's method for optimization;
    \item distinguish Newton's method for optimization from Newton root
          finding;
    \item recognize modified Newton methods;
    \item understand damped Newton methods;
    \item understand quasi-Newton methods;
    \item understand the secant condition;
    \item recognize the BFGS update;
    \item recognize DFP conceptually;
    \item understand limited-memory BFGS;
    \item understand conjugate-gradient optimization conceptually;
    \item understand nonlinear conjugate gradient;
    \item recognize trust-region methods;
    \item formulate the trust-region subproblem;
    \item understand predicted and actual reduction;
    \item understand trust-region radius updates;
    \item compare line-search and trust-region globalization;
    \item understand derivative-free optimization conceptually;
    \item recognize coordinate-search and simplex-type methods;
    \item understand finite-difference gradients;
    \item understand automatic differentiation conceptually;
    \item distinguish forward and reverse automatic differentiation;
    \item understand constrained numerical optimization;
    \item understand projected-gradient methods;
    \item recognize penalty methods;
    \item recognize barrier methods;
    \item understand augmented Lagrangian methods;
    \item understand sequential quadratic programming;
    \item understand interior-point methods conceptually;
    \item recognize primal-dual systems;
    \item understand active-set methods;
    \item understand numerical KKT residuals;
    \item formulate stopping criteria;
    \item understand scaling and preconditioning in optimization;
    \item recognize sparse optimization structure;
    \item understand large-scale optimization;
    \item recognize stochastic gradient methods;
    \item distinguish batch, stochastic, and mini-batch gradients;
    \item understand momentum conceptually;
    \item understand adaptive first-order methods conceptually;
    \item recognize proximal-gradient methods;
    \item understand nonsmooth optimization conceptually;
    \item connect numerical optimization with machine learning, statistics,
          engineering, optimal control, inverse problems, and scientific
          computing.
\end{itemize}
}

%%%%%%%%%%%%%%%%%%%%%%%%%%%%%%%%%%%%%%%%%%%%%%%%%%%%%%%%%%%%
\subsection{What Is Numerical Optimization?}
\label{subsec:what-is-numerical-optimization}
%%%%%%%%%%%%%%%%%%%%%%%%%%%%%%%%%%%%%%%%%%%%%%%%%%%%%%%%%%%%

Numerical optimization develops algorithms for approximately solving
optimization problems using finite computation.

A general problem has the form

\[
\boxed{
\min_{x\in\mathbb{R}^n}
f(x)
}
\]

or, with constraints,

\[
\boxed{
\begin{aligned}
\min_x
\quad&
f(x)
\\
\text{subject to}
\quad&
g_i(x)\leq0,
\\
&
h_j(x)=0.
\end{aligned}
}
\]

The goal is not merely to describe optimality conditions.

The goal is to construct a computational procedure that finds a useful
solution.

%%%%%%%%%%%%%%%%%%%%%%%%%%%%%%%%%%%%%%%%%%%%%%%%%%%%%%%%%%%%
\subsection{Iterative Optimization}
\label{subsec:iterative-optimization}
%%%%%%%%%%%%%%%%%%%%%%%%%%%%%%%%%%%%%%%%%%%%%%%%%%%%%%%%%%%%

Most nonlinear numerical optimization algorithms generate iterates

\[
\boxed{
x_0,x_1,x_2,\ldots.
}
\]

A common update is

\[
\boxed{
x_{k+1}
=
x_k
+
\alpha_kp_k,
}
\]

where

\[
p_k
=
\text{search direction},
\]

and

\[
\alpha_k
=
\text{step length}.
\]

The algorithm must decide both quantities effectively.

%%%%%%%%%%%%%%%%%%%%%%%%%%%%%%%%%%%%%%%%%%%%%%%%%%%%%%%%%%%%
\subsection{Descent Directions}
\label{subsec:descent-directions}
%%%%%%%%%%%%%%%%%%%%%%%%%%%%%%%%%%%%%%%%%%%%%%%%%%%%%%%%%%%%

Suppose \(f\) is differentiable.

A direction \(p\) is a descent direction at \(x\) if

\[
\boxed{
\nabla f(x)^Tp<0.
}
\]

Then for sufficiently small positive \(\alpha\),

\[
\boxed{
f(x+\alpha p)<f(x).
}
\]

This follows from the directional derivative.

%%%%%%%%%%%%%%%%%%%%%%%%%%%%%%%%%%%%%%%%%%%%%%%%%%%%%%%%%%%%
\subsection{The Negative Gradient}
\label{subsec:negative-gradient-direction}
%%%%%%%%%%%%%%%%%%%%%%%%%%%%%%%%%%%%%%%%%%%%%%%%%%%%%%%%%%%%

The negative gradient is

\[
\boxed{
p
=
-\nabla f(x).
}
\]

Then

\[
\nabla f(x)^Tp
=
-
\|\nabla f(x)\|_2^2.
\]

If

\[
\nabla f(x)\neq0,
\]

then

\[
\boxed{
\nabla f(x)^Tp<0.
}
\]

Thus the negative gradient is always a descent direction at a
nonstationary point.

%%%%%%%%%%%%%%%%%%%%%%%%%%%%%%%%%%%%%%%%%%%%%%%%%%%%%%%%%%%%
\subsection{Gradient Descent}
\label{subsec:gradient-descent}
%%%%%%%%%%%%%%%%%%%%%%%%%%%%%%%%%%%%%%%%%%%%%%%%%%%%%%%%%%%%

Gradient descent uses

\[
\boxed{
x_{k+1}
=
x_k
-
\alpha_k
\nabla f(x_k).
}
\]

The main computational question is how to choose

\[
\alpha_k.
\]

Possible choices include:

\[
\boxed{
\text{constant step size},
}
\]

\[
\boxed{
\text{exact line search},
}
\]

or

\[
\boxed{
\text{inexact line search}.
}
\]

%%%%%%%%%%%%%%%%%%%%%%%%%%%%%%%%%%%%%%%%%%%%%%%%%%%%%%%%%%%%
\subsection{Line Search}
\label{subsec:numerical-line-search}
%%%%%%%%%%%%%%%%%%%%%%%%%%%%%%%%%%%%%%%%%%%%%%%%%%%%%%%%%%%%

Given search direction \(p_k\), define the one-dimensional function

\[
\boxed{
\phi(\alpha)
=
f(x_k+\alpha p_k).
}
\]

A line search seeks a useful value of

\[
\alpha>0.
\]

The full multidimensional problem is temporarily reduced to a
one-dimensional search along \(p_k\).

%%%%%%%%%%%%%%%%%%%%%%%%%%%%%%%%%%%%%%%%%%%%%%%%%%%%%%%%%%%%
\subsection{Exact Line Search}
\label{subsec:exact-line-search}
%%%%%%%%%%%%%%%%%%%%%%%%%%%%%%%%%%%%%%%%%%%%%%%%%%%%%%%%%%%%

An exact line search chooses

\[
\boxed{
\alpha_k
\in
\operatorname*{arg\,min}_{\alpha>0}
f(x_k+\alpha p_k).
}
\]

This minimizes the objective exactly along the chosen line.

In practice, exact line search can require too much computation relative
to the benefit gained.

%%%%%%%%%%%%%%%%%%%%%%%%%%%%%%%%%%%%%%%%%%%%%%%%%%%%%%%%%%%%
\subsection{Inexact Line Search}
\label{subsec:inexact-line-search}
%%%%%%%%%%%%%%%%%%%%%%%%%%%%%%%%%%%%%%%%%%%%%%%%%%%%%%%%%%%%

An inexact line search does not find the exact minimizer of

\[
\phi(\alpha).
\]

Instead, it seeks a step satisfying conditions that guarantee sufficient
progress.

Common conditions include:

\[
\boxed{
\text{Armijo},
}
\]

\[
\boxed{
\text{Wolfe},
}
\]

and

\[
\boxed{
\text{strong Wolfe}.
}
\]

%%%%%%%%%%%%%%%%%%%%%%%%%%%%%%%%%%%%%%%%%%%%%%%%%%%%%%%%%%%%
\subsection{Armijo Condition}
\label{subsec:armijo-condition}
%%%%%%%%%%%%%%%%%%%%%%%%%%%%%%%%%%%%%%%%%%%%%%%%%%%%%%%%%%%%

The Armijo sufficient-decrease condition is

\[
\boxed{
f(x_k+\alpha p_k)
\leq
f(x_k)
+
c_1
\alpha
\nabla f(x_k)^Tp_k,
}
\]

where

\[
\boxed{
0<c_1<1.
}
\]

Because

\[
\nabla f(x_k)^Tp_k<0,
\]

the right-hand side is smaller than \(f(x_k)\).

%%%%%%%%%%%%%%%%%%%%%%%%%%%%%%%%%%%%%%%%%%%%%%%%%%%%%%%%%%%%
\subsection{Backtracking Line Search}
\label{subsec:backtracking-line-search}
%%%%%%%%%%%%%%%%%%%%%%%%%%%%%%%%%%%%%%%%%%%%%%%%%%%%%%%%%%%%

A common Armijo procedure is:

\begin{enumerate}
    \item choose initial trial step \(\alpha>0\);
    \item test the Armijo condition;
    \item if it fails, replace
    \[
    \alpha
    \leftarrow
    \rho\alpha,
    \]
    where
    \[
    0<\rho<1;
    \]
    \item repeat until sufficient decrease holds.
\end{enumerate}

This is called backtracking line search.

%%%%%%%%%%%%%%%%%%%%%%%%%%%%%%%%%%%%%%%%%%%%%%%%%%%%%%%%%%%%
\subsection{Worked Example: Backtracking}
\label{subsec:worked-backtracking}
%%%%%%%%%%%%%%%%%%%%%%%%%%%%%%%%%%%%%%%%%%%%%%%%%%%%%%%%%%%%

\begin{workedexamplebox}{Testing an Armijo Step}

Let

\[
f(x)=x^2,
\]

with

\[
x_k=2.
\]

Then

\[
\nabla f(2)=4.
\]

Choose steepest-descent direction

\[
p_k=-4.
\]

Take

\[
c_1=0.1
\]

and trial step

\[
\alpha=0.5.
\]

The candidate point is

\[
x_{k+1}
=
2+0.5(-4)
=
0.
\]

Thus

\[
f(x_{k+1})=0.
\]

The Armijo right-hand side is

\[
4
+
0.1(0.5)(4)(-4)
=
4-0.8
=
3.2.
\]

Since

\[
0\leq3.2,
\]

the condition is satisfied.

\begin{answerbox}
\[
\boxed{
\alpha=0.5
\text{ is accepted by this Armijo test.}
}
\]
\end{answerbox}

\end{workedexamplebox}

%%%%%%%%%%%%%%%%%%%%%%%%%%%%%%%%%%%%%%%%%%%%%%%%%%%%%%%%%%%%
\subsection{Wolfe Conditions}
\label{subsec:wolfe-conditions}
%%%%%%%%%%%%%%%%%%%%%%%%%%%%%%%%%%%%%%%%%%%%%%%%%%%%%%%%%%%%

The Wolfe conditions combine sufficient decrease with a curvature
condition.

The first condition is Armijo:

\[
f(x_k+\alpha p_k)
\leq
f(x_k)
+
c_1\alpha
\nabla f(x_k)^Tp_k.
\]

The curvature condition is

\[
\boxed{
\nabla f(x_k+\alpha p_k)^Tp_k
\geq
c_2
\nabla f(x_k)^Tp_k,
}
\]

where

\[
\boxed{
0<c_1<c_2<1.
}
\]

%%%%%%%%%%%%%%%%%%%%%%%%%%%%%%%%%%%%%%%%%%%%%%%%%%%%%%%%%%%%
\subsection{Strong Wolfe Conditions}
\label{subsec:strong-wolfe}
%%%%%%%%%%%%%%%%%%%%%%%%%%%%%%%%%%%%%%%%%%%%%%%%%%%%%%%%%%%%

The strong Wolfe curvature condition is

\[
\boxed{
\left|
\nabla f(x_k+\alpha p_k)^Tp_k
\right|
\leq
c_2
\left|
\nabla f(x_k)^Tp_k
\right|.
}
\]

This prevents the step from stopping too early while also controlling the
magnitude of the slope after the step.

Strong Wolfe conditions are especially important in nonlinear conjugate
gradient and quasi-Newton methods.

%%%%%%%%%%%%%%%%%%%%%%%%%%%%%%%%%%%%%%%%%%%%%%%%%%%%%%%%%%%%
\subsection{Quadratic Objective}
\label{subsec:quadratic-objective-numerical-optimization}
%%%%%%%%%%%%%%%%%%%%%%%%%%%%%%%%%%%%%%%%%%%%%%%%%%%%%%%%%%%%

Consider

\[
\boxed{
f(x)
=
\frac12x^TAx-b^Tx,
}
\]

where

\[
A=A^T\succ0.
\]

Then

\[
\boxed{
\nabla f(x)
=
Ax-b,
}
\]

and

\[
\boxed{
\nabla^2f(x)
=
A.
}
\]

The unique minimizer satisfies

\[
\boxed{
Ax^*=b.
}
\]

Thus minimizing a positive-definite quadratic is equivalent to solving a
linear system.

%%%%%%%%%%%%%%%%%%%%%%%%%%%%%%%%%%%%%%%%%%%%%%%%%%%%%%%%%%%%
\subsection{Exact Step for Steepest Descent on a Quadratic}
\label{subsec:quadratic-steepest-descent-step}
%%%%%%%%%%%%%%%%%%%%%%%%%%%%%%%%%%%%%%%%%%%%%%%%%%%%%%%%%%%%

Let

\[
g_k
=
\nabla f(x_k)
=
Ax_k-b.
\]

Choose

\[
p_k=-g_k.
\]

For exact line search,

\[
\boxed{
\alpha_k
=
\frac{
g_k^Tg_k
}{
g_k^TAg_k
}.
}
\]

This gives the best step along the negative-gradient direction.

%%%%%%%%%%%%%%%%%%%%%%%%%%%%%%%%%%%%%%%%%%%%%%%%%%%%%%%%%%%%
\subsection{Conditioning and Gradient Descent}
\label{subsec:conditioning-gradient-descent}
%%%%%%%%%%%%%%%%%%%%%%%%%%%%%%%%%%%%%%%%%%%%%%%%%%%%%%%%%%%%

For positive-definite quadratic objectives, gradient descent can converge
slowly when

\[
\boxed{
\kappa_2(A)
}
\]

is large.

Geometrically, level sets become highly elongated ellipses.

The negative gradient may repeatedly zigzag across a narrow valley.

Thus poor conditioning slows first-order optimization.

%%%%%%%%%%%%%%%%%%%%%%%%%%%%%%%%%%%%%%%%%%%%%%%%%%%%%%%%%%%%
\subsection{Worked Example: Gradient Descent Step}
\label{subsec:worked-gradient-descent-step}
%%%%%%%%%%%%%%%%%%%%%%%%%%%%%%%%%%%%%%%%%%%%%%%%%%%%%%%%%%%%

\begin{workedexamplebox}{One Gradient Step}

Let

\[
f(x,y)
=
x^2+2y^2.
\]

Then

\[
\nabla f(x,y)
=
\begin{pmatrix}
2x\\
4y
\end{pmatrix}.
\]

At

\[
(x_0,y_0)
=
(1,1),
\]

\[
\nabla f(1,1)
=
\begin{pmatrix}
2\\
4
\end{pmatrix}.
\]

With step size

\[
\alpha=0.1,
\]

gradient descent gives

\[
\begin{pmatrix}
x_1\\
y_1
\end{pmatrix}
=
\begin{pmatrix}
1\\
1
\end{pmatrix}
-
0.1
\begin{pmatrix}
2\\
4
\end{pmatrix}.
\]

Thus,

\begin{answerbox}
\[
\boxed{
(x_1,y_1)
=
(0.8,0.6).
}
\]
\end{answerbox}

\end{workedexamplebox}

%%%%%%%%%%%%%%%%%%%%%%%%%%%%%%%%%%%%%%%%%%%%%%%%%%%%%%%%%%%%
\subsection{Newton's Method for Optimization}
\label{subsec:newton-method-optimization}
%%%%%%%%%%%%%%%%%%%%%%%%%%%%%%%%%%%%%%%%%%%%%%%%%%%%%%%%%%%%

Near \(x_k\), approximate \(f\) by its second-order Taylor model:

\[
\boxed{
m_k(p)
=
f(x_k)
+
g_k^Tp
+
\frac12
p^TH_kp,
}
\]

where

\[
g_k
=
\nabla f(x_k),
\]

and

\[
H_k
=
\nabla^2f(x_k).
\]

Minimizing the quadratic model gives

\[
\boxed{
H_kp_k
=
-g_k.
}
\]

%%%%%%%%%%%%%%%%%%%%%%%%%%%%%%%%%%%%%%%%%%%%%%%%%%%%%%%%%%%%
\subsection{Newton Step}
\label{subsec:newton-step-optimization}
%%%%%%%%%%%%%%%%%%%%%%%%%%%%%%%%%%%%%%%%%%%%%%%%%%%%%%%%%%%%

If \(H_k\) is nonsingular,

\[
\boxed{
p_k
=
-
H_k^{-1}g_k.
}
\]

The update is

\[
\boxed{
x_{k+1}
=
x_k+p_k.
}
\]

In implementation, one solves

\[
H_kp_k=-g_k
\]

rather than explicitly computing \(H_k^{-1}\).

%%%%%%%%%%%%%%%%%%%%%%%%%%%%%%%%%%%%%%%%%%%%%%%%%%%%%%%%%%%%
\subsection{Newton Optimization Versus Newton Root Finding}
\label{subsec:newton-optimization-vs-root}
%%%%%%%%%%%%%%%%%%%%%%%%%%%%%%%%%%%%%%%%%%%%%%%%%%%%%%%%%%%%

A local minimizer satisfies

\[
\boxed{
\nabla f(x)=0.
}
\]

Applying Newton root finding to the vector equation

\[
\nabla f(x)=0
\]

produces

\[
\boxed{
x_{k+1}
=
x_k
-
[
\nabla^2f(x_k)
]^{-1}
\nabla f(x_k).
}
\]

Thus Newton optimization can be viewed as Newton's method applied to the
first-order optimality equation.

%%%%%%%%%%%%%%%%%%%%%%%%%%%%%%%%%%%%%%%%%%%%%%%%%%%%%%%%%%%%
\subsection{Quadratic Convergence of Newton's Method}
\label{subsec:newton-optimization-quadratic}
%%%%%%%%%%%%%%%%%%%%%%%%%%%%%%%%%%%%%%%%%%%%%%%%%%%%%%%%%%%%

Near a nondegenerate local minimizer \(x^*\), under suitable smoothness
conditions and with

\[
\nabla^2f(x^*)\succ0,
\]

Newton's method can converge quadratically:

\[
\boxed{
\|x_{k+1}-x^*\|
\leq
C
\|x_k-x^*\|^2.
}
\]

This rapid local convergence is one of Newton's major advantages.

%%%%%%%%%%%%%%%%%%%%%%%%%%%%%%%%%%%%%%%%%%%%%%%%%%%%%%%%%%%%
\subsection{When Newton's Direction Is a Descent Direction}
\label{subsec:newton-descent-direction}
%%%%%%%%%%%%%%%%%%%%%%%%%%%%%%%%%%%%%%%%%%%%%%%%%%%%%%%%%%%%

If

\[
H_k\succ0,
\]

then

\[
p_k
=
-H_k^{-1}g_k.
\]

Therefore,

\[
g_k^Tp_k
=
-
g_k^TH_k^{-1}g_k.
\]

For \(g_k\neq0\),

\[
\boxed{
g_k^Tp_k<0.
}
\]

Thus the Newton direction is a descent direction when the Hessian is
positive definite.

%%%%%%%%%%%%%%%%%%%%%%%%%%%%%%%%%%%%%%%%%%%%%%%%%%%%%%%%%%%%
\subsection{Indefinite Hessians}
\label{subsec:indefinite-hessians}
%%%%%%%%%%%%%%%%%%%%%%%%%%%%%%%%%%%%%%%%%%%%%%%%%%%%%%%%%%%%

If

\[
H_k
\]

is indefinite, the Newton direction may fail to be a descent direction.

This can occur far from a minimizer or near saddle points.

Practical algorithms therefore modify or globalize Newton's method.

%%%%%%%%%%%%%%%%%%%%%%%%%%%%%%%%%%%%%%%%%%%%%%%%%%%%%%%%%%%%
\subsection{Damped Newton Method}
\label{subsec:damped-newton}
%%%%%%%%%%%%%%%%%%%%%%%%%%%%%%%%%%%%%%%%%%%%%%%%%%%%%%%%%%%%

Instead of always taking the full Newton step, use

\[
\boxed{
x_{k+1}
=
x_k
+
\alpha_kp_k,
}
\]

where \(p_k\) solves

\[
H_kp_k=-g_k.
\]

A line search chooses

\[
0<\alpha_k\leq1.
\]

Far from the solution, damping improves robustness.

Near the solution, the full step

\[
\alpha_k=1
\]

is often eventually accepted.

%%%%%%%%%%%%%%%%%%%%%%%%%%%%%%%%%%%%%%%%%%%%%%%%%%%%%%%%%%%%
\subsection{Modified Newton Method}
\label{subsec:modified-newton}
%%%%%%%%%%%%%%%%%%%%%%%%%%%%%%%%%%%%%%%%%%%%%%%%%%%%%%%%%%%%

If the Hessian is not positive definite, replace it by a modified matrix

\[
\boxed{
B_k
\succ0.
}
\]

For example,

\[
\boxed{
B_k
=
H_k+\tau_kI
}
\]

with sufficiently large

\[
\tau_k>0.
\]

Then solve

\[
\boxed{
B_kp_k=-g_k.
}
\]

This forces a descent direction.

%%%%%%%%%%%%%%%%%%%%%%%%%%%%%%%%%%%%%%%%%%%%%%%%%%%%%%%%%%%%
\subsection{Quasi-Newton Methods}
\label{subsec:quasi-newton-methods}
%%%%%%%%%%%%%%%%%%%%%%%%%%%%%%%%%%%%%%%%%%%%%%%%%%%%%%%%%%%%

Quasi-Newton methods approximate Hessian information using gradient
changes instead of explicitly computing second derivatives.

Define

\[
\boxed{
s_k
=
x_{k+1}-x_k,
}
\]

and

\[
\boxed{
y_k
=
\nabla f(x_{k+1})
-
\nabla f(x_k).
}
\]

These quantities approximate the action of the Hessian.

%%%%%%%%%%%%%%%%%%%%%%%%%%%%%%%%%%%%%%%%%%%%%%%%%%%%%%%%%%%%
\subsection{Secant Condition}
\label{subsec:quasi-newton-secant-condition}
%%%%%%%%%%%%%%%%%%%%%%%%%%%%%%%%%%%%%%%%%%%%%%%%%%%%%%%%%%%%

A Hessian approximation \(B_{k+1}\) is often required to satisfy

\[
\boxed{
B_{k+1}s_k
=
y_k.
}
\]

This is the quasi-Newton secant condition.

It is the multidimensional analogue of replacing derivatives by secant
information.

%%%%%%%%%%%%%%%%%%%%%%%%%%%%%%%%%%%%%%%%%%%%%%%%%%%%%%%%%%%%
\subsection{Inverse Secant Condition}
\label{subsec:inverse-secant-condition}
%%%%%%%%%%%%%%%%%%%%%%%%%%%%%%%%%%%%%%%%%%%%%%%%%%%%%%%%%%%%

If \(H_k\) denotes an approximation to the inverse Hessian, then the
corresponding condition is

\[
\boxed{
H_{k+1}y_k
=
s_k.
}
\]

The search direction is then

\[
\boxed{
p_k
=
-H_kg_k.
}
\]

%%%%%%%%%%%%%%%%%%%%%%%%%%%%%%%%%%%%%%%%%%%%%%%%%%%%%%%%%%%%
\subsection{BFGS Update}
\label{subsec:bfgs-update}
%%%%%%%%%%%%%%%%%%%%%%%%%%%%%%%%%%%%%%%%%%%%%%%%%%%%%%%%%%%%

One of the most widely used quasi-Newton methods is BFGS.

For a Hessian approximation \(B_k\), the update is

\[
\boxed{
B_{k+1}
=
B_k
-
\frac{
B_ks_ks_k^TB_k
}{
s_k^TB_ks_k
}
+
\frac{
y_ky_k^T
}{
y_k^Ts_k
}.
}
\]

Under

\[
\boxed{
y_k^Ts_k>0,
}
\]

positive definiteness is preserved when \(B_k\succ0\).

%%%%%%%%%%%%%%%%%%%%%%%%%%%%%%%%%%%%%%%%%%%%%%%%%%%%%%%%%%%%
\subsection{Inverse BFGS Update}
\label{subsec:inverse-bfgs-update}
%%%%%%%%%%%%%%%%%%%%%%%%%%%%%%%%%%%%%%%%%%%%%%%%%%%%%%%%%%%%

Let

\[
\boxed{
\rho_k
=
\frac1{y_k^Ts_k}.
}
\]

The inverse-Hessian update can be written

\[
\boxed{
H_{k+1}
=
(
I-\rho_ks_ky_k^T
)
H_k
(
I-\rho_ky_ks_k^T
)
+
\rho_ks_ks_k^T.
}
\]

Then

\[
\boxed{
p_k
=
-H_kg_k.
}
\]

%%%%%%%%%%%%%%%%%%%%%%%%%%%%%%%%%%%%%%%%%%%%%%%%%%%%%%%%%%%%
\subsection{Why BFGS Is Important}
\label{subsec:why-bfgs}
%%%%%%%%%%%%%%%%%%%%%%%%%%%%%%%%%%%%%%%%%%%%%%%%%%%%%%%%%%%%

BFGS combines:

\[
\boxed{
\text{first-derivative information},
}
\]

with an evolving approximation to

\[
\boxed{
\text{second-order curvature}.
}
\]

It often achieves much faster convergence than ordinary gradient descent
without requiring explicit Hessian evaluation.

%%%%%%%%%%%%%%%%%%%%%%%%%%%%%%%%%%%%%%%%%%%%%%%%%%%%%%%%%%%%
\subsection{DFP Method Preview}
\label{subsec:dfp-preview}
%%%%%%%%%%%%%%%%%%%%%%%%%%%%%%%%%%%%%%%%%%%%%%%%%%%%%%%%%%%%

The Davidon--Fletcher--Powell update is another classical quasi-Newton
method.

Like BFGS, it satisfies a secant equation and updates curvature
information.

BFGS is usually preferred in many modern unconstrained applications
because of its practical numerical behavior.

%%%%%%%%%%%%%%%%%%%%%%%%%%%%%%%%%%%%%%%%%%%%%%%%%%%%%%%%%%%%
\subsection{Limited-Memory BFGS}
\label{subsec:limited-memory-bfgs}
%%%%%%%%%%%%%%%%%%%%%%%%%%%%%%%%%%%%%%%%%%%%%%%%%%%%%%%%%%%%

For very large \(n\), storing an \(n\times n\) Hessian approximation is
expensive.

Limited-memory BFGS, abbreviated L-BFGS, stores only a small number of
recent pairs

\[
\boxed{
(s_k,y_k).
}
\]

This reduces storage from approximately

\[
O(n^2)
\]

to approximately

\[
\boxed{
O(mn),
}
\]

where \(m\) is the number of stored correction pairs.

%%%%%%%%%%%%%%%%%%%%%%%%%%%%%%%%%%%%%%%%%%%%%%%%%%%%%%%%%%%%
\subsection{Nonlinear Conjugate Gradient}
\label{subsec:nonlinear-conjugate-gradient}
%%%%%%%%%%%%%%%%%%%%%%%%%%%%%%%%%%%%%%%%%%%%%%%%%%%%%%%%%%%%

Nonlinear conjugate-gradient methods generate directions

\[
\boxed{
p_{k+1}
=
-g_{k+1}
+
\beta_kp_k.
}
\]

The coefficient \(\beta_k\) is selected using gradient information.

Different formulas produce different nonlinear conjugate-gradient
variants.

%%%%%%%%%%%%%%%%%%%%%%%%%%%%%%%%%%%%%%%%%%%%%%%%%%%%%%%%%%%%
\subsection{Fletcher--Reeves Parameter}
\label{subsec:fletcher-reeves}
%%%%%%%%%%%%%%%%%%%%%%%%%%%%%%%%%%%%%%%%%%%%%%%%%%%%%%%%%%%%

One common choice is

\[
\boxed{
\beta_k^{FR}
=
\frac{
g_{k+1}^Tg_{k+1}
}{
g_k^Tg_k
}.
}
\]

Combined with appropriate line searches, this extends the conjugate-
gradient idea beyond quadratic functions.

%%%%%%%%%%%%%%%%%%%%%%%%%%%%%%%%%%%%%%%%%%%%%%%%%%%%%%%%%%%%
\subsection{Polak--Ribiere Parameter}
\label{subsec:polak-ribiere}
%%%%%%%%%%%%%%%%%%%%%%%%%%%%%%%%%%%%%%%%%%%%%%%%%%%%%%%%%%%%

Another common choice is

\[
\boxed{
\beta_k^{PR}
=
\frac{
g_{k+1}^T
(
g_{k+1}-g_k
)
}{
g_k^Tg_k
}.
}
\]

Practical implementations may use safeguarded variants such as

\[
\boxed{
\max
\{
\beta_k^{PR},0
\}.
}
\]

%%%%%%%%%%%%%%%%%%%%%%%%%%%%%%%%%%%%%%%%%%%%%%%%%%%%%%%%%%%%
\subsection{Trust-Region Methods}
\label{subsec:trust-region-methods}
%%%%%%%%%%%%%%%%%%%%%%%%%%%%%%%%%%%%%%%%%%%%%%%%%%%%%%%%%%%%

A trust-region method constructs a local model

\[
m_k(p)
\]

of \(f(x_k+p)\).

Instead of trusting the model everywhere, constrain the step:

\[
\boxed{
\|p\|
\leq
\Delta_k.
}
\]

The quantity

\[
\Delta_k
\]

is the trust-region radius.

%%%%%%%%%%%%%%%%%%%%%%%%%%%%%%%%%%%%%%%%%%%%%%%%%%%%%%%%%%%%
\subsection{Quadratic Trust-Region Model}
\label{subsec:quadratic-trust-region-model}
%%%%%%%%%%%%%%%%%%%%%%%%%%%%%%%%%%%%%%%%%%%%%%%%%%%%%%%%%%%%

A standard model is

\[
\boxed{
m_k(p)
=
f_k
+
g_k^Tp
+
\frac12p^TB_kp.
}
\]

The trust-region subproblem is

\[
\boxed{
\min_{\|p\|\leq\Delta_k}
m_k(p).
}
\]

The matrix \(B_k\) may be the exact Hessian or an approximation.

%%%%%%%%%%%%%%%%%%%%%%%%%%%%%%%%%%%%%%%%%%%%%%%%%%%%%%%%%%%%
\subsection{Actual and Predicted Reduction}
\label{subsec:actual-predicted-reduction}
%%%%%%%%%%%%%%%%%%%%%%%%%%%%%%%%%%%%%%%%%%%%%%%%%%%%%%%%%%%%

The actual objective reduction is

\[
\boxed{
f(x_k)-f(x_k+p_k).
}
\]

The model predicts reduction

\[
\boxed{
m_k(0)-m_k(p_k).
}
\]

Their ratio is

\[
\boxed{
\rho_k
=
\frac{
f(x_k)-f(x_k+p_k)
}{
m_k(0)-m_k(p_k)
}.
}
\]

%%%%%%%%%%%%%%%%%%%%%%%%%%%%%%%%%%%%%%%%%%%%%%%%%%%%%%%%%%%%
\subsection{Trust-Region Acceptance}
\label{subsec:trust-region-acceptance}
%%%%%%%%%%%%%%%%%%%%%%%%%%%%%%%%%%%%%%%%%%%%%%%%%%%%%%%%%%%%

If

\[
\rho_k
\]

is large and positive, the model predicted the actual behavior well.

The step is typically accepted.

If

\[
\rho_k
\]

is poor or negative, the model was unreliable over the proposed region.

The step may be rejected and

\[
\Delta_k
\]

reduced.

%%%%%%%%%%%%%%%%%%%%%%%%%%%%%%%%%%%%%%%%%%%%%%%%%%%%%%%%%%%%
\subsection{Updating the Trust Radius}
\label{subsec:trust-radius-update}
%%%%%%%%%%%%%%%%%%%%%%%%%%%%%%%%%%%%%%%%%%%%%%%%%%%%%%%%%%%%

A typical rule is:

\[
\boxed{
\rho_k\text{ poor}
\Rightarrow
\Delta_{k+1}<\Delta_k,
}
\]

\[
\boxed{
\rho_k\text{ good}
\Rightarrow
\Delta_{k+1}\approx\Delta_k,
}
\]

and when the model is very accurate and the boundary is reached,

\[
\boxed{
\Delta_{k+1}>\Delta_k.
}
\]

Exact numerical thresholds vary by implementation.

%%%%%%%%%%%%%%%%%%%%%%%%%%%%%%%%%%%%%%%%%%%%%%%%%%%%%%%%%%%%
\subsection{Cauchy Point}
\label{subsec:trust-region-cauchy-point}
%%%%%%%%%%%%%%%%%%%%%%%%%%%%%%%%%%%%%%%%%%%%%%%%%%%%%%%%%%%%

A simple trust-region step minimizes the quadratic model along the
negative-gradient direction:

\[
p=-\tau g_k,
\]

subject to

\[
\|p\|\leq\Delta_k.
\]

The best such point is called the Cauchy point.

It provides a guaranteed amount of model decrease under standard
assumptions.

%%%%%%%%%%%%%%%%%%%%%%%%%%%%%%%%%%%%%%%%%%%%%%%%%%%%%%%%%%%%
\subsection{Line Search Versus Trust Region}
\label{subsec:line-search-vs-trust-region}
%%%%%%%%%%%%%%%%%%%%%%%%%%%%%%%%%%%%%%%%%%%%%%%%%%%%%%%%%%%%

Line-search methods first choose a direction and then determine how far to
move.

Trust-region methods first determine how far the local model should be
trusted and then optimize within that region.

Thus:

\[
\boxed{
\text{line search}
=
\text{direction first},
}
\]

while

\[
\boxed{
\text{trust region}
=
\text{region first}.
}
\]

Both provide globalization beyond purely local Newton steps.

%%%%%%%%%%%%%%%%%%%%%%%%%%%%%%%%%%%%%%%%%%%%%%%%%%%%%%%%%%%%
\subsection{Derivative-Free Optimization}
\label{subsec:derivative-free-optimization}
%%%%%%%%%%%%%%%%%%%%%%%%%%%%%%%%%%%%%%%%%%%%%%%%%%%%%%%%%%%%

Some optimization problems do not provide reliable derivatives.

Examples include:

\[
\boxed{
\text{black-box simulations},
}
\]

\[
\boxed{
\text{noisy experiments},
}
\]

or

\[
\boxed{
\text{nonsmooth computational models}.
}
\]

Derivative-free optimization uses objective evaluations without requiring
analytic gradients.

%%%%%%%%%%%%%%%%%%%%%%%%%%%%%%%%%%%%%%%%%%%%%%%%%%%%%%%%%%%%
\subsection{Coordinate Search}
\label{subsec:coordinate-search}
%%%%%%%%%%%%%%%%%%%%%%%%%%%%%%%%%%%%%%%%%%%%%%%%%%%%%%%%%%%%

A simple derivative-free method tests movement along coordinate
directions:

\[
\boxed{
\pm e_1,
\pm e_2,
\ldots,
\pm e_n.
}
\]

If a direction improves the objective, the algorithm moves there.

Otherwise the search step may be reduced.

%%%%%%%%%%%%%%%%%%%%%%%%%%%%%%%%%%%%%%%%%%%%%%%%%%%%%%%%%%%%
\subsection{Simplex-Type Direct Search}
\label{subsec:simplex-direct-search}
%%%%%%%%%%%%%%%%%%%%%%%%%%%%%%%%%%%%%%%%%%%%%%%%%%%%%%%%%%%%

Simplex-type direct-search methods maintain a collection of

\[
n+1
\]

points in \(n\)-dimensional space.

The geometry is modified using operations such as:

\[
\boxed{
\text{reflection},
}
\]

\[
\boxed{
\text{expansion},
}
\]

\[
\boxed{
\text{contraction},
}
\]

and

\[
\boxed{
\text{shrinkage}.
}
\]

Such methods can be useful for small derivative-free problems.

%%%%%%%%%%%%%%%%%%%%%%%%%%%%%%%%%%%%%%%%%%%%%%%%%%%%%%%%%%%%
\subsection{Finite-Difference Gradients}
\label{subsec:finite-difference-gradients}
%%%%%%%%%%%%%%%%%%%%%%%%%%%%%%%%%%%%%%%%%%%%%%%%%%%%%%%%%%%%

If analytic derivatives are unavailable, approximate

\[
\frac{\partial f}{\partial x_i}
\]

using

\[
\boxed{
\frac{
f(x+he_i)-f(x)
}{
h
}.
}
\]

This requires approximately one additional objective evaluation per
variable for a forward-difference gradient.

%%%%%%%%%%%%%%%%%%%%%%%%%%%%%%%%%%%%%%%%%%%%%%%%%%%%%%%%%%%%
\subsection{Centered Gradient Approximation}
\label{subsec:centered-gradient-approximation}
%%%%%%%%%%%%%%%%%%%%%%%%%%%%%%%%%%%%%%%%%%%%%%%%%%%%%%%%%%%%

A centered approximation is

\[
\boxed{
\frac{
\partial f
}{
\partial x_i
}
\approx
\frac{
f(x+he_i)
-
f(x-he_i)
}{
2h
}.
}
\]

This has higher truncation accuracy but requires approximately twice as
many evaluations.

%%%%%%%%%%%%%%%%%%%%%%%%%%%%%%%%%%%%%%%%%%%%%%%%%%%%%%%%%%%%
\subsection{Finite-Difference Step Selection}
\label{subsec:optimization-finite-difference-step}
%%%%%%%%%%%%%%%%%%%%%%%%%%%%%%%%%%%%%%%%%%%%%%%%%%%%%%%%%%%%

A very large \(h\) causes truncation error.

A very small \(h\) amplifies floating-point cancellation.

Therefore derivative approximation requires a compromise between:

\[
\boxed{
\text{truncation error}
}
\]

and

\[
\boxed{
\text{roundoff error}.
}
\]

%%%%%%%%%%%%%%%%%%%%%%%%%%%%%%%%%%%%%%%%%%%%%%%%%%%%%%%%%%%%
\subsection{Automatic Differentiation}
\label{subsec:automatic-differentiation}
%%%%%%%%%%%%%%%%%%%%%%%%%%%%%%%%%%%%%%%%%%%%%%%%%%%%%%%%%%%%

Automatic differentiation, abbreviated AD, computes derivatives by
systematically applying the chain rule to the elementary operations used
to evaluate a function.

It is not the same as:

\[
\boxed{
\text{symbolic differentiation},
}
\]

and it is not the same as:

\[
\boxed{
\text{finite differences}.
}
\]

AD can often provide derivatives accurate to floating-point precision.

%%%%%%%%%%%%%%%%%%%%%%%%%%%%%%%%%%%%%%%%%%%%%%%%%%%%%%%%%%%%
\subsection{Forward-Mode Automatic Differentiation}
\label{subsec:forward-mode-ad}
%%%%%%%%%%%%%%%%%%%%%%%%%%%%%%%%%%%%%%%%%%%%%%%%%%%%%%%%%%%%

Forward-mode AD propagates:

\[
\boxed{
\text{function values}
}
\]

and

\[
\boxed{
\text{directional derivatives}
}
\]

through the computational graph.

It is especially effective when the number of inputs is small relative to
the number of outputs.

%%%%%%%%%%%%%%%%%%%%%%%%%%%%%%%%%%%%%%%%%%%%%%%%%%%%%%%%%%%%
\subsection{Reverse-Mode Automatic Differentiation}
\label{subsec:reverse-mode-ad}
%%%%%%%%%%%%%%%%%%%%%%%%%%%%%%%%%%%%%%%%%%%%%%%%%%%%%%%%%%%%

Reverse-mode AD first evaluates the function and records the computational
graph.

Then derivative information is propagated backward.

For a scalar objective

\[
f:\mathbb{R}^n\to\mathbb{R},
\]

reverse mode can compute the entire gradient at a cost proportional to a
small multiple of the original function evaluation.

This makes it fundamental in large-scale machine learning.

%%%%%%%%%%%%%%%%%%%%%%%%%%%%%%%%%%%%%%%%%%%%%%%%%%%%%%%%%%%%
\subsection{Gradient Checking}
\label{subsec:gradient-checking}
%%%%%%%%%%%%%%%%%%%%%%%%%%%%%%%%%%%%%%%%%%%%%%%%%%%%%%%%%%%%

Analytic or automatic gradients can be checked numerically by comparing
them with finite-difference directional derivatives.

For direction \(d\),

\[
\boxed{
\nabla f(x)^Td
\approx
\frac{
f(x+hd)-f(x)
}{
h
}.
}
\]

Agreement across an appropriate range of \(h\) helps detect derivative
implementation errors.

%%%%%%%%%%%%%%%%%%%%%%%%%%%%%%%%%%%%%%%%%%%%%%%%%%%%%%%%%%%%
\subsection{Constrained Numerical Optimization}
\label{subsec:constrained-numerical-optimization}
%%%%%%%%%%%%%%%%%%%%%%%%%%%%%%%%%%%%%%%%%%%%%%%%%%%%%%%%%%%%

Consider

\[
\boxed{
\begin{aligned}
\min_x
\quad&
f(x)
\\
\text{subject to}
\quad&
g_i(x)\leq0,
\\
&
h_j(x)=0.
\end{aligned}
}
\]

Numerical methods must improve the objective while maintaining or
restoring feasibility.

This makes constrained optimization substantially more complicated than
unconstrained optimization.

%%%%%%%%%%%%%%%%%%%%%%%%%%%%%%%%%%%%%%%%%%%%%%%%%%%%%%%%%%%%
\subsection{Projected Gradient Method}
\label{subsec:projected-gradient-method}
%%%%%%%%%%%%%%%%%%%%%%%%%%%%%%%%%%%%%%%%%%%%%%%%%%%%%%%%%%%%

Suppose \(C\) is a closed feasible set and projection is computationally
tractable.

Projected gradient uses

\[
\boxed{
x_{k+1}
=
P_C
\left(
x_k
-
\alpha_k
\nabla f(x_k)
\right),
}
\]

where

\[
\boxed{
P_C(y)
=
\operatorname*{arg\,min}_{x\in C}
\|x-y\|_2.
}
\]

For convex \(C\), the Euclidean projection is uniquely defined.

%%%%%%%%%%%%%%%%%%%%%%%%%%%%%%%%%%%%%%%%%%%%%%%%%%%%%%%%%%%%
\subsection{Bound Constraints}
\label{subsec:bound-constrained-optimization}
%%%%%%%%%%%%%%%%%%%%%%%%%%%%%%%%%%%%%%%%%%%%%%%%%%%%%%%%%%%%

A common feasible set is

\[
\boxed{
\ell_i
\leq
x_i
\leq
u_i.
}
\]

Projection is then componentwise:

\[
\boxed{
[P_C(y)]_i
=
\min
\left\{
u_i,
\max
\{
\ell_i,y_i
\}
\right\}.
}
\]

This makes projected methods simple for box constraints.

%%%%%%%%%%%%%%%%%%%%%%%%%%%%%%%%%%%%%%%%%%%%%%%%%%%%%%%%%%%%
\subsection{Penalty Methods}
\label{subsec:numerical-penalty-methods}
%%%%%%%%%%%%%%%%%%%%%%%%%%%%%%%%%%%%%%%%%%%%%%%%%%%%%%%%%%%%

Penalty methods replace constrained optimization by a sequence of
unconstrained problems.

For equality constraints

\[
h(x)=0,
\]

a quadratic penalty objective is

\[
\boxed{
\Phi_\mu(x)
=
f(x)
+
\frac{\mu}{2}
\|h(x)\|_2^2.
}
\]

As

\[
\mu
\]

increases, constraint violations become increasingly expensive.

%%%%%%%%%%%%%%%%%%%%%%%%%%%%%%%%%%%%%%%%%%%%%%%%%%%%%%%%%%%%
\subsection{Penalty Ill Conditioning}
\label{subsec:penalty-ill-conditioning}
%%%%%%%%%%%%%%%%%%%%%%%%%%%%%%%%%%%%%%%%%%%%%%%%%%%%%%%%%%%%

Very large penalty parameters can force near-feasibility.

However, they may also make the objective badly conditioned.

Thus a pure penalty method can become numerically difficult as

\[
\mu\to\infty.
\]

This motivates augmented Lagrangian approaches.

%%%%%%%%%%%%%%%%%%%%%%%%%%%%%%%%%%%%%%%%%%%%%%%%%%%%%%%%%%%%
\subsection{Barrier Methods}
\label{subsec:numerical-barrier-methods}
%%%%%%%%%%%%%%%%%%%%%%%%%%%%%%%%%%%%%%%%%%%%%%%%%%%%%%%%%%%%

For inequality constraints

\[
g_i(x)<0,
\]

a logarithmic barrier is

\[
\boxed{
\Phi_\mu(x)
=
f(x)
-
\mu
\sum_i
\ln
(
-g_i(x)
),
}
\]

where

\[
\mu>0.
\]

The barrier tends to \(+\infty\) as a constraint boundary is approached
from the feasible interior.

%%%%%%%%%%%%%%%%%%%%%%%%%%%%%%%%%%%%%%%%%%%%%%%%%%%%%%%%%%%%
\subsection{Central Path}
\label{subsec:numerical-central-path}
%%%%%%%%%%%%%%%%%%%%%%%%%%%%%%%%%%%%%%%%%%%%%%%%%%%%%%%%%%%%

For each barrier parameter \(\mu\), solve the barrier subproblem.

The collection of minimizers

\[
\boxed{
x(\mu)
}
\]

forms a central-path-type trajectory.

As

\[
\mu\to0,
\]

the barrier solution approaches the constrained optimum under suitable
assumptions.

%%%%%%%%%%%%%%%%%%%%%%%%%%%%%%%%%%%%%%%%%%%%%%%%%%%%%%%%%%%%
\subsection{Augmented Lagrangian}
\label{subsec:numerical-augmented-lagrangian}
%%%%%%%%%%%%%%%%%%%%%%%%%%%%%%%%%%%%%%%%%%%%%%%%%%%%%%%%%%%%

For equality constraints

\[
h(x)=0,
\]

the augmented Lagrangian is

\[
\boxed{
\mathcal{L}_\mu(x,\lambda)
=
f(x)
+
\lambda^Th(x)
+
\frac{\mu}{2}
\|h(x)\|_2^2.
}
\]

This combines multiplier information with a penalty term.

%%%%%%%%%%%%%%%%%%%%%%%%%%%%%%%%%%%%%%%%%%%%%%%%%%%%%%%%%%%%
\subsection{Multiplier Update}
\label{subsec:augmented-lagrangian-multiplier-update}
%%%%%%%%%%%%%%%%%%%%%%%%%%%%%%%%%%%%%%%%%%%%%%%%%%%%%%%%%%%%

A basic multiplier update is

\[
\boxed{
\lambda_{k+1}
=
\lambda_k
+
\mu_k
h(x_{k+1}).
}
\]

The multiplier corrects systematic constraint error, reducing the need for
extremely large penalty parameters.

%%%%%%%%%%%%%%%%%%%%%%%%%%%%%%%%%%%%%%%%%%%%%%%%%%%%%%%%%%%%
\subsection{Sequential Quadratic Programming}
\label{subsec:numerical-sqp}
%%%%%%%%%%%%%%%%%%%%%%%%%%%%%%%%%%%%%%%%%%%%%%%%%%%%%%%%%%%%

Sequential Quadratic Programming, abbreviated SQP, solves a sequence of
quadratic approximations to a nonlinear constrained problem.

At \(x_k\), linearize constraints and approximate the Lagrangian Hessian.

A typical SQP subproblem is

\[
\boxed{
\begin{aligned}
\min_p
\quad&
\nabla f(x_k)^Tp
+
\frac12p^TB_kp
\\
\text{subject to}
\quad&
g_i(x_k)
+
\nabla g_i(x_k)^Tp
\leq0,
\\
&
h_j(x_k)
+
\nabla h_j(x_k)^Tp
=
0.
\end{aligned}
}
\]

%%%%%%%%%%%%%%%%%%%%%%%%%%%%%%%%%%%%%%%%%%%%%%%%%%%%%%%%%%%%
\subsection{SQP Interpretation}
\label{subsec:sqp-interpretation}
%%%%%%%%%%%%%%%%%%%%%%%%%%%%%%%%%%%%%%%%%%%%%%%%%%%%%%%%%%%%

SQP can be viewed as the constrained analogue of Newton's method.

Newton solves a local quadratic model of the objective.

SQP solves a local quadratic model while also imposing linearized
constraints.

Near a regular solution, SQP can converge very rapidly.

%%%%%%%%%%%%%%%%%%%%%%%%%%%%%%%%%%%%%%%%%%%%%%%%%%%%%%%%%%%%
\subsection{Merit Functions}
\label{subsec:optimization-merit-functions}
%%%%%%%%%%%%%%%%%%%%%%%%%%%%%%%%%%%%%%%%%%%%%%%%%%%%%%%%%%%%

A constrained method may reduce the objective while worsening feasibility,
or improve feasibility while increasing the objective.

A merit function combines both effects.

One example is

\[
\boxed{
\phi(x)
=
f(x)
+
\rho
\|h(x)\|_1
+
\rho
\|[g(x)]_+\|_1.
}
\]

The merit function helps decide whether a trial step makes overall
progress.

%%%%%%%%%%%%%%%%%%%%%%%%%%%%%%%%%%%%%%%%%%%%%%%%%%%%%%%%%%%%
\subsection{Filter Methods}
\label{subsec:optimization-filter-methods}
%%%%%%%%%%%%%%%%%%%%%%%%%%%%%%%%%%%%%%%%%%%%%%%%%%%%%%%%%%%%

A filter method avoids combining feasibility and objective values into one
weighted scalar.

Instead, it accepts steps that sufficiently improve either:

\[
\boxed{
\text{constraint violation}
}
\]

or

\[
\boxed{
\text{objective value}.
}
\]

This can reduce sensitivity to merit-function penalty parameters.

%%%%%%%%%%%%%%%%%%%%%%%%%%%%%%%%%%%%%%%%%%%%%%%%%%%%%%%%%%%%
\subsection{Active-Set Methods}
\label{subsec:numerical-active-set-methods}
%%%%%%%%%%%%%%%%%%%%%%%%%%%%%%%%%%%%%%%%%%%%%%%%%%%%%%%%%%%%

At a constrained solution, some inequalities are active:

\[
\boxed{
g_i(x)=0.
}
\]

Active-set methods estimate which constraints should be active and solve
equality-constrained subproblems based on that working set.

The working set changes as the algorithm proceeds.

%%%%%%%%%%%%%%%%%%%%%%%%%%%%%%%%%%%%%%%%%%%%%%%%%%%%%%%%%%%%
\subsection{Interior-Point Methods}
\label{subsec:numerical-interior-point-methods}
%%%%%%%%%%%%%%%%%%%%%%%%%%%%%%%%%%%%%%%%%%%%%%%%%%%%%%%%%%%%

Interior-point methods approach the optimum through interior feasible
points or approximate interior iterates.

They are closely connected with barrier functions and primal-dual KKT
systems.

Modern interior-point methods are central to:

\[
\boxed{
\text{linear programming},
}
\]

\[
\boxed{
\text{quadratic programming},
}
\]

and

\[
\boxed{
\text{nonlinear optimization}.
}
\]

%%%%%%%%%%%%%%%%%%%%%%%%%%%%%%%%%%%%%%%%%%%%%%%%%%%%%%%%%%%%
\subsection{Primal-Dual Interior-Point Idea}
\label{subsec:primal-dual-interior-point}
%%%%%%%%%%%%%%%%%%%%%%%%%%%%%%%%%%%%%%%%%%%%%%%%%%%%%%%%%%%%

For inequality constraint

\[
g_i(x)\leq0,
\]

with multiplier

\[
\lambda_i\geq0,
\]

ordinary complementarity is

\[
\boxed{
\lambda_i
g_i(x)
=
0.
}
\]

Interior-point methods replace exact complementarity with a perturbed
relation, conceptually of the form

\[
\boxed{
\lambda_i
(
-g_i(x)
)
\approx
\mu.
}
\]

As

\[
\mu\to0,
\]

the iterates approach complementarity.

%%%%%%%%%%%%%%%%%%%%%%%%%%%%%%%%%%%%%%%%%%%%%%%%%%%%%%%%%%%%
\subsection{KKT Residuals}
\label{subsec:numerical-kkt-residuals}
%%%%%%%%%%%%%%%%%%%%%%%%%%%%%%%%%%%%%%%%%%%%%%%%%%%%%%%%%%%%

A numerical constrained solver should monitor first-order optimality
conditions.

Typical residuals include:

\[
\boxed{
r_{\mathrm{stationarity}}
=
\nabla_x
\mathcal{L}(x,\lambda,\nu),
}
\]

\[
\boxed{
r_{\mathrm{eq}}
=
h(x),
}
\]

and complementarity residuals such as

\[
\boxed{
r_{\mathrm{comp},i}
=
\lambda_i g_i(x).
}
\]

Constraint violation must also be monitored.

%%%%%%%%%%%%%%%%%%%%%%%%%%%%%%%%%%%%%%%%%%%%%%%%%%%%%%%%%%%%
\subsection{Stopping Criteria}
\label{subsec:optimization-stopping-criteria}
%%%%%%%%%%%%%%%%%%%%%%%%%%%%%%%%%%%%%%%%%%%%%%%%%%%%%%%%%%%%

A numerical optimization algorithm may stop when several measures are
small.

Examples include:

\[
\boxed{
\|\nabla f(x_k)\|
\leq
\varepsilon_g,
}
\]

\[
\boxed{
\|x_{k+1}-x_k\|
\leq
\varepsilon_x,
}
\]

and

\[
\boxed{
|f(x_{k+1})-f(x_k)|
\leq
\varepsilon_f.
}
\]

For constrained problems, feasibility and KKT residuals must also be
checked.

%%%%%%%%%%%%%%%%%%%%%%%%%%%%%%%%%%%%%%%%%%%%%%%%%%%%%%%%%%%%
\subsection{Relative Stopping Criteria}
\label{subsec:relative-optimization-stopping}
%%%%%%%%%%%%%%%%%%%%%%%%%%%%%%%%%%%%%%%%%%%%%%%%%%%%%%%%%%%%

A relative step test may use

\[
\boxed{
\frac{
\|x_{k+1}-x_k\|
}{
1+\|x_k\|
}
\leq
\varepsilon.
}
\]

A relative objective test may use

\[
\boxed{
\frac{
|f_{k+1}-f_k|
}{
1+|f_k|
}
\leq
\varepsilon.
}
\]

Relative tests account for problem scale.

%%%%%%%%%%%%%%%%%%%%%%%%%%%%%%%%%%%%%%%%%%%%%%%%%%%%%%%%%%%%
\subsection{Stopping Too Early}
\label{subsec:optimization-stopping-too-early}
%%%%%%%%%%%%%%%%%%%%%%%%%%%%%%%%%%%%%%%%%%%%%%%%%%%%%%%%%%%%

A small change in objective value does not necessarily imply stationarity.

The algorithm may simply be taking very small steps.

Similarly, a small gradient does not automatically imply a minimum when the
Hessian is indefinite.

Multiple diagnostics should be used.

%%%%%%%%%%%%%%%%%%%%%%%%%%%%%%%%%%%%%%%%%%%%%%%%%%%%%%%%%%%%
\subsection{Scaling}
\label{subsec:numerical-optimization-scaling}
%%%%%%%%%%%%%%%%%%%%%%%%%%%%%%%%%%%%%%%%%%%%%%%%%%%%%%%%%%%%

Optimization variables may have dramatically different magnitudes.

For example,

\[
x_1
\approx10^{-6},
\qquad
x_2
\approx10^6.
\]

Poor scaling can distort:

\[
\boxed{
\text{gradients},
}
\]

\[
\boxed{
\text{Hessians},
}
\]

and

\[
\boxed{
\text{step lengths}.
}
\]

Scaling can substantially improve numerical performance.

%%%%%%%%%%%%%%%%%%%%%%%%%%%%%%%%%%%%%%%%%%%%%%%%%%%%%%%%%%%%
\subsection{Variable Scaling}
\label{subsec:variable-scaling}
%%%%%%%%%%%%%%%%%%%%%%%%%%%%%%%%%%%%%%%%%%%%%%%%%%%%%%%%%%%%

Introduce scaled variables

\[
\boxed{
x=Dz,
}
\]

where \(D\) is diagonal.

Then solve in \(z\)-coordinates.

The goal is to make typical components of \(z\) have comparable
magnitudes.

%%%%%%%%%%%%%%%%%%%%%%%%%%%%%%%%%%%%%%%%%%%%%%%%%%%%%%%%%%%%
\subsection{Objective Scaling}
\label{subsec:objective-scaling}
%%%%%%%%%%%%%%%%%%%%%%%%%%%%%%%%%%%%%%%%%%%%%%%%%%%%%%%%%%%%

Multiplying an objective by a positive constant does not change the exact
minimizer:

\[
\boxed{
\operatorname*{arg\,min}_x f(x)
=
\operatorname*{arg\,min}_x cf(x),
\qquad
c>0.
}
\]

However, the numerical magnitude of derivatives and tolerances changes.

This can affect solver performance.

%%%%%%%%%%%%%%%%%%%%%%%%%%%%%%%%%%%%%%%%%%%%%%%%%%%%%%%%%%%%
\subsection{Constraint Scaling}
\label{subsec:constraint-scaling}
%%%%%%%%%%%%%%%%%%%%%%%%%%%%%%%%%%%%%%%%%%%%%%%%%%%%%%%%%%%%

Suppose one constraint has typical magnitude

\[
10^{-8}
\]

while another has magnitude

\[
10^6.
\]

Treating their raw residuals equally can be numerically misleading.

Constraint scaling makes feasibility tests more comparable across
constraints.

%%%%%%%%%%%%%%%%%%%%%%%%%%%%%%%%%%%%%%%%%%%%%%%%%%%%%%%%%%%%
\subsection{Preconditioning in Optimization}
\label{subsec:optimization-preconditioning}
%%%%%%%%%%%%%%%%%%%%%%%%%%%%%%%%%%%%%%%%%%%%%%%%%%%%%%%%%%%%

Gradient descent uses the Euclidean geometry of the coordinates.

A preconditioned step uses

\[
\boxed{
p_k
=
-
M_k^{-1}g_k,
}
\]

where

\[
M_k\succ0.
\]

If \(M_k\) approximates the Hessian well, this can reduce effects of poor
conditioning.

Newton and quasi-Newton methods can be viewed as sophisticated
preconditioned gradient methods.

%%%%%%%%%%%%%%%%%%%%%%%%%%%%%%%%%%%%%%%%%%%%%%%%%%%%%%%%%%%%
\subsection{Sparse Optimization}
\label{subsec:sparse-numerical-optimization}
%%%%%%%%%%%%%%%%%%%%%%%%%%%%%%%%%%%%%%%%%%%%%%%%%%%%%%%%%%%%

Large optimization problems often contain sparse derivatives.

Examples include:

\[
\boxed{
\text{sparse Jacobians},
}
\]

\[
\boxed{
\text{sparse Hessians},
}
\]

and

\[
\boxed{
\text{sparse constraint matrices}.
}
\]

Numerical methods should exploit this structure rather than storing dense
matrices.

%%%%%%%%%%%%%%%%%%%%%%%%%%%%%%%%%%%%%%%%%%%%%%%%%%%%%%%%%%%%
\subsection{Newton--KKT Systems}
\label{subsec:newton-kkt-systems}
%%%%%%%%%%%%%%%%%%%%%%%%%%%%%%%%%%%%%%%%%%%%%%%%%%%%%%%%%%%%

Equality-constrained Newton or SQP methods may require solving systems of
the form

\[
\boxed{
\begin{pmatrix}
H & A^T\\
A & 0
\end{pmatrix}
\begin{pmatrix}
p\\
\lambda
\end{pmatrix}
=
-
\begin{pmatrix}
g\\
c
\end{pmatrix}.
}
\]

These are saddle-point or KKT systems.

Efficient numerical linear algebra is essential for solving them.

%%%%%%%%%%%%%%%%%%%%%%%%%%%%%%%%%%%%%%%%%%%%%%%%%%%%%%%%%%%%
\subsection{Schur Complement Methods}
\label{subsec:optimization-schur-complement}
%%%%%%%%%%%%%%%%%%%%%%%%%%%%%%%%%%%%%%%%%%%%%%%%%%%%%%%%%%%%

Block linear systems can sometimes be reduced using a Schur complement.

If \(H\) is invertible, eliminate \(p\) to obtain a smaller multiplier
system involving

\[
\boxed{
AH^{-1}A^T.
}
\]

This can reduce computational cost when the constraint dimension is much
smaller than the variable dimension.

%%%%%%%%%%%%%%%%%%%%%%%%%%%%%%%%%%%%%%%%%%%%%%%%%%%%%%%%%%%%
\subsection{Hessian-Vector Products}
\label{subsec:hessian-vector-products}
%%%%%%%%%%%%%%%%%%%%%%%%%%%%%%%%%%%%%%%%%%%%%%%%%%%%%%%%%%%%

Large-scale algorithms may not form the full Hessian.

Instead, they compute products

\[
\boxed{
\nabla^2f(x)v.
}
\]

Such products can be obtained using automatic differentiation or
directional derivative techniques.

This supports matrix-free Newton--Krylov methods.

%%%%%%%%%%%%%%%%%%%%%%%%%%%%%%%%%%%%%%%%%%%%%%%%%%%%%%%%%%%%
\subsection{Newton--CG Methods}
\label{subsec:newton-cg-methods}
%%%%%%%%%%%%%%%%%%%%%%%%%%%%%%%%%%%%%%%%%%%%%%%%%%%%%%%%%%%%

If the Hessian is positive definite or appropriately modified, the Newton
system

\[
\boxed{
H_kp_k=-g_k
}
\]

can be solved approximately using conjugate gradient.

The linear system need not be solved to high precision when far from the
optimum.

This can greatly reduce large-scale computational cost.

%%%%%%%%%%%%%%%%%%%%%%%%%%%%%%%%%%%%%%%%%%%%%%%%%%%%%%%%%%%%
\subsection{Inexact Newton Methods}
\label{subsec:inexact-newton-methods}
%%%%%%%%%%%%%%%%%%%%%%%%%%%%%%%%%%%%%%%%%%%%%%%%%%%%%%%%%%%%

An inexact Newton step satisfies

\[
\boxed{
H_kp_k
=
-g_k+r_k,
}
\]

where \(r_k\) is a controlled residual.

Early iterations may use loose linear-solve tolerances.

As the optimization approaches the solution, the linear-system tolerance
can be tightened.

%%%%%%%%%%%%%%%%%%%%%%%%%%%%%%%%%%%%%%%%%%%%%%%%%%%%%%%%%%%%
\subsection{Nonsmooth Optimization}
\label{subsec:nonsmooth-numerical-optimization}
%%%%%%%%%%%%%%%%%%%%%%%%%%%%%%%%%%%%%%%%%%%%%%%%%%%%%%%%%%%%

Some important objectives are not differentiable everywhere.

Examples include:

\[
\boxed{
|x|,
}
\]

\[
\boxed{
\|x\|_1,
}
\]

and

\[
\boxed{
\max_i f_i(x).
}
\]

Ordinary Newton and gradient methods may not apply directly at
nondifferentiable points.

%%%%%%%%%%%%%%%%%%%%%%%%%%%%%%%%%%%%%%%%%%%%%%%%%%%%%%%%%%%%
\subsection{Subgradient Method}
\label{subsec:subgradient-method}
%%%%%%%%%%%%%%%%%%%%%%%%%%%%%%%%%%%%%%%%%%%%%%%%%%%%%%%%%%%%

For convex nonsmooth \(f\), choose

\[
g_k
\in
\partial f(x_k),
\]

where \(\partial f(x_k)\) is the subdifferential.

Then iterate

\[
\boxed{
x_{k+1}
=
x_k
-
\alpha_kg_k.
}
\]

Subgradient methods are simple but usually converge more slowly than smooth
gradient methods.

%%%%%%%%%%%%%%%%%%%%%%%%%%%%%%%%%%%%%%%%%%%%%%%%%%%%%%%%%%%%
\subsection{Composite Optimization}
\label{subsec:composite-optimization}
%%%%%%%%%%%%%%%%%%%%%%%%%%%%%%%%%%%%%%%%%%%%%%%%%%%%%%%%%%%%

A common problem is

\[
\boxed{
\min_x
f(x)+g(x),
}
\]

where:

\[
f
=
\text{smooth function},
\]

and

\[
g
=
\text{possibly nonsmooth but structured function}.
\]

This structure leads to proximal methods.

%%%%%%%%%%%%%%%%%%%%%%%%%%%%%%%%%%%%%%%%%%%%%%%%%%%%%%%%%%%%
\subsection{Proximal Operator}
\label{subsec:proximal-operator-numerical}
%%%%%%%%%%%%%%%%%%%%%%%%%%%%%%%%%%%%%%%%%%%%%%%%%%%%%%%%%%%%

The proximal operator of \(g\) is

\[
\boxed{
\operatorname{prox}_{\alpha g}(v)
=
\operatorname*{arg\,min}_x
\left\{
g(x)
+
\frac1{2\alpha}
\|x-v\|_2^2
\right\}.
}
\]

It generalizes projection and provides a way to handle structured
nonsmooth terms.

%%%%%%%%%%%%%%%%%%%%%%%%%%%%%%%%%%%%%%%%%%%%%%%%%%%%%%%%%%%%
\subsection{Proximal Gradient Method}
\label{subsec:proximal-gradient-method}
%%%%%%%%%%%%%%%%%%%%%%%%%%%%%%%%%%%%%%%%%%%%%%%%%%%%%%%%%%%%

For

\[
f(x)+g(x),
\]

with smooth \(f\), proximal gradient uses

\[
\boxed{
x_{k+1}
=
\operatorname{prox}_{\alpha_kg}
\left(
x_k
-
\alpha_k
\nabla f(x_k)
\right).
}
\]

This combines a gradient step on the smooth part with a proximal step on
the nonsmooth part.

%%%%%%%%%%%%%%%%%%%%%%%%%%%%%%%%%%%%%%%%%%%%%%%%%%%%%%%%%%%%
\subsection{Soft Thresholding}
\label{subsec:soft-thresholding-numerical}
%%%%%%%%%%%%%%%%%%%%%%%%%%%%%%%%%%%%%%%%%%%%%%%%%%%%%%%%%%%%

For

\[
g(x)
=
\lambda\|x\|_1,
\]

the proximal operator is componentwise soft thresholding:

\[
\boxed{
[\operatorname{prox}_{\alpha\lambda\|\cdot\|_1}(v)]_i
=
\operatorname{sign}(v_i)
\max
\{
|v_i|-\alpha\lambda,
0
\}.
}
\]

This is central to sparse optimization.

%%%%%%%%%%%%%%%%%%%%%%%%%%%%%%%%%%%%%%%%%%%%%%%%%%%%%%%%%%%%
\subsection{Stochastic Gradient Descent}
\label{subsec:stochastic-gradient-descent}
%%%%%%%%%%%%%%%%%%%%%%%%%%%%%%%%%%%%%%%%%%%%%%%%%%%%%%%%%%%%

Suppose

\[
\boxed{
f(x)
=
\frac1N
\sum_{i=1}^{N}
f_i(x).
}
\]

Full gradient descent uses

\[
\boxed{
\nabla f(x)
=
\frac1N
\sum_{i=1}^{N}
\nabla f_i(x).
}
\]

For very large \(N\), this can be expensive.

Stochastic Gradient Descent, abbreviated SGD, uses only one or a small
subset of terms at each step.

%%%%%%%%%%%%%%%%%%%%%%%%%%%%%%%%%%%%%%%%%%%%%%%%%%%%%%%%%%%%
\subsection{Single-Sample SGD}
\label{subsec:single-sample-sgd}
%%%%%%%%%%%%%%%%%%%%%%%%%%%%%%%%%%%%%%%%%%%%%%%%%%%%%%%%%%%%

Choose random index \(i_k\).

Then

\[
\boxed{
x_{k+1}
=
x_k
-
\alpha_k
\nabla f_{i_k}(x_k).
}
\]

If the sampling is appropriate,

\[
\boxed{
\mathbb{E}
[
\nabla f_{i_k}(x)
]
=
\nabla f(x).
}
\]

The stochastic gradient is then an unbiased estimator of the full
gradient.

%%%%%%%%%%%%%%%%%%%%%%%%%%%%%%%%%%%%%%%%%%%%%%%%%%%%%%%%%%%%
\subsection{Mini-Batch Gradient}
\label{subsec:mini-batch-gradient}
%%%%%%%%%%%%%%%%%%%%%%%%%%%%%%%%%%%%%%%%%%%%%%%%%%%%%%%%%%%%

Let \(B_k\) be a random subset of data indices.

Define

\[
\boxed{
g_k
=
\frac1{|B_k|}
\sum_{i\in B_k}
\nabla f_i(x_k).
}
\]

Then update

\[
\boxed{
x_{k+1}
=
x_k-\alpha_kg_k.
}
\]

Mini-batches balance gradient variance against computational efficiency.

%%%%%%%%%%%%%%%%%%%%%%%%%%%%%%%%%%%%%%%%%%%%%%%%%%%%%%%%%%%%
\subsection{Batch, Mini-Batch, and Stochastic Gradients}
\label{subsec:batch-mini-stochastic}
%%%%%%%%%%%%%%%%%%%%%%%%%%%%%%%%%%%%%%%%%%%%%%%%%%%%%%%%%%%%

Full batch:

\[
\boxed{
|B_k|=N.
}
\]

Pure stochastic:

\[
\boxed{
|B_k|=1.
}
\]

Mini-batch:

\[
\boxed{
1<|B_k|<N.
}
\]

The appropriate choice depends on computational architecture, noise, and
problem scale.

%%%%%%%%%%%%%%%%%%%%%%%%%%%%%%%%%%%%%%%%%%%%%%%%%%%%%%%%%%%%
\subsection{Stochastic Step Sizes}
\label{subsec:stochastic-step-sizes}
%%%%%%%%%%%%%%%%%%%%%%%%%%%%%%%%%%%%%%%%%%%%%%%%%%%%%%%%%%%%

Classical stochastic approximation often uses decreasing step sizes
satisfying conditions such as

\[
\boxed{
\sum_{k=0}^{\infty}
\alpha_k
=
\infty,
}
\]

and

\[
\boxed{
\sum_{k=0}^{\infty}
\alpha_k^2
<
\infty.
}
\]

These conditions balance continued movement with diminishing stochastic
noise under suitable assumptions.

%%%%%%%%%%%%%%%%%%%%%%%%%%%%%%%%%%%%%%%%%%%%%%%%%%%%%%%%%%%%
\subsection{Constant Learning Rates}
\label{subsec:constant-learning-rates}
%%%%%%%%%%%%%%%%%%%%%%%%%%%%%%%%%%%%%%%%%%%%%%%%%%%%%%%%%%%%

In practical large-scale learning problems, constant or piecewise-constant
step sizes are common.

A constant step can move rapidly but may cause the iterates to fluctuate
around an optimum rather than converge exactly.

Reducing the step size later can improve final accuracy.

%%%%%%%%%%%%%%%%%%%%%%%%%%%%%%%%%%%%%%%%%%%%%%%%%%%%%%%%%%%%
\subsection{Momentum}
\label{subsec:optimization-momentum}
%%%%%%%%%%%%%%%%%%%%%%%%%%%%%%%%%%%%%%%%%%%%%%%%%%%%%%%%%%%%

Momentum introduces a velocity-like variable.

A basic form is

\[
\boxed{
v_{k+1}
=
\beta v_k
+
g_k,
}
\]

\[
\boxed{
x_{k+1}
=
x_k
-
\alpha v_{k+1},
}
\]

where

\[
0\leq\beta<1.
\]

Momentum can smooth oscillations and accelerate movement along persistent
descent directions.

%%%%%%%%%%%%%%%%%%%%%%%%%%%%%%%%%%%%%%%%%%%%%%%%%%%%%%%%%%%%
\subsection{Accelerated Gradient Preview}
\label{subsec:accelerated-gradient-preview}
%%%%%%%%%%%%%%%%%%%%%%%%%%%%%%%%%%%%%%%%%%%%%%%%%%%%%%%%%%%%

For smooth convex optimization, acceleration methods combine gradient
steps with carefully designed momentum terms.

Under standard assumptions, accelerated methods can improve the
first-order convergence rate compared with ordinary gradient descent.

The analysis requires more structure than simply adding arbitrary
momentum.

%%%%%%%%%%%%%%%%%%%%%%%%%%%%%%%%%%%%%%%%%%%%%%%%%%%%%%%%%%%%
\subsection{Adaptive First-Order Methods Preview}
\label{subsec:adaptive-first-order-preview}
%%%%%%%%%%%%%%%%%%%%%%%%%%%%%%%%%%%%%%%%%%%%%%%%%%%%%%%%%%%%

Adaptive first-order methods modify step sizes using past gradient
information.

They may maintain running estimates of:

\[
\boxed{
\text{gradient magnitudes},
}
\]

or

\[
\boxed{
\text{first and second moments}.
}
\]

These methods are widely used in machine learning.

Their practical behavior depends strongly on hyperparameter choices and
problem structure.

%%%%%%%%%%%%%%%%%%%%%%%%%%%%%%%%%%%%%%%%%%%%%%%%%%%%%%%%%%%%
\subsection{Deterministic Versus Stochastic Optimization Algorithms}
\label{subsec:deterministic-versus-stochastic-algorithms}
%%%%%%%%%%%%%%%%%%%%%%%%%%%%%%%%%%%%%%%%%%%%%%%%%%%%%%%%%%%%

Deterministic algorithms evaluate an exact objective or gradient for the
current model.

Stochastic algorithms use random approximations.

Thus deterministic iterations are often more predictable, while
stochastic iterations may be much cheaper.

Large data sets frequently favor stochastic or mini-batch methods.

%%%%%%%%%%%%%%%%%%%%%%%%%%%%%%%%%%%%%%%%%%%%%%%%%%%%%%%%%%%%
\subsection{Local and Global Optimization}
\label{subsec:local-global-numerical-optimization}
%%%%%%%%%%%%%%%%%%%%%%%%%%%%%%%%%%%%%%%%%%%%%%%%%%%%%%%%%%%%

Most numerical nonlinear optimization algorithms are designed to find a
local stationary point.

For nonconvex \(f\), the result may depend strongly on the initial guess.

Finding a global optimum may require:

\[
\boxed{
\text{multiple starts},
}
\]

\[
\boxed{
\text{branch-and-bound},
}
\]

\[
\boxed{
\text{global relaxations},
}
\]

or specialized global-search methods.

%%%%%%%%%%%%%%%%%%%%%%%%%%%%%%%%%%%%%%%%%%%%%%%%%%%%%%%%%%%%
\subsection{Saddle Points}
\label{subsec:optimization-saddle-points}
%%%%%%%%%%%%%%%%%%%%%%%%%%%%%%%%%%%%%%%%%%%%%%%%%%%%%%%%%%%%

A point satisfying

\[
\boxed{
\nabla f(x)=0
}
\]

need not be a minimizer.

If the Hessian has both positive and negative eigenvalues, the point is a
saddle point.

Second-order information can help distinguish minima from saddles.

%%%%%%%%%%%%%%%%%%%%%%%%%%%%%%%%%%%%%%%%%%%%%%%%%%%%%%%%%%%%
\subsection{Negative Curvature}
\label{subsec:negative-curvature}
%%%%%%%%%%%%%%%%%%%%%%%%%%%%%%%%%%%%%%%%%%%%%%%%%%%%%%%%%%%%

If there exists \(d\) such that

\[
\boxed{
d^THd<0,
}
\]

then \(d\) is a direction of negative curvature.

Near a stationary point, moving along such a direction can decrease the
objective.

Trust-region and modified Newton methods can exploit negative curvature
rather than blindly applying a Newton step.

%%%%%%%%%%%%%%%%%%%%%%%%%%%%%%%%%%%%%%%%%%%%%%%%%%%%%%%%%%%%
\subsection{Optimization Landscapes}
\label{subsec:optimization-landscapes}
%%%%%%%%%%%%%%%%%%%%%%%%%%%%%%%%%%%%%%%%%%%%%%%%%%%%%%%%%%%%

A nonconvex objective may contain:

\[
\boxed{
\text{multiple local minima},
}
\]

\[
\boxed{
\text{local maxima},
}
\]

\[
\boxed{
\text{saddle points},
}
\]

and

\[
\boxed{
\text{flat regions}.
}
\]

Numerical algorithms interact with this geometry differently.

Initialization and globalization become important.

%%%%%%%%%%%%%%%%%%%%%%%%%%%%%%%%%%%%%%%%%%%%%%%%%%%%%%%%%%%%
\subsection{Warm Starts}
\label{subsec:optimization-warm-starts}
%%%%%%%%%%%%%%%%%%%%%%%%%%%%%%%%%%%%%%%%%%%%%%%%%%%%%%%%%%%%

A warm start initializes a solver using a solution or approximate solution
from a related optimization problem.

This is useful when solving sequences of nearby problems.

Examples include:

\[
\boxed{
\text{MPC},
}
\]

\[
\boxed{
\text{parameter continuation},
}
\]

and

\[
\boxed{
\text{regularization paths}.
}
\]

%%%%%%%%%%%%%%%%%%%%%%%%%%%%%%%%%%%%%%%%%%%%%%%%%%%%%%%%%%%%
\subsection{Continuation Methods}
\label{subsec:optimization-continuation}
%%%%%%%%%%%%%%%%%%%%%%%%%%%%%%%%%%%%%%%%%%%%%%%%%%%%%%%%%%%%

Suppose a difficult problem depends on parameter \(\mu\).

Start with an easier value

\[
\mu_0,
\]

solve that problem, then gradually change \(\mu\).

Each previous solution becomes the initial guess for the next problem.

Barrier methods and homotopy methods use this principle.

%%%%%%%%%%%%%%%%%%%%%%%%%%%%%%%%%%%%%%%%%%%%%%%%%%%%%%%%%%%%
\subsection{Numerical Optimization Workflow}
\label{subsec:numerical-optimization-workflow}
%%%%%%%%%%%%%%%%%%%%%%%%%%%%%%%%%%%%%%%%%%%%%%%%%%%%%%%%%%%%

A practical workflow is:

\begin{enumerate}
    \item formulate the objective and constraints;
    \item determine whether the problem is smooth;
    \item identify convexity or nonconvexity;
    \item determine whether gradients are available;
    \item determine whether Hessians or Hessian-vector products are
          available;
    \item scale variables, objectives, and constraints;
    \item choose an initial point;
    \item determine whether feasibility is required initially;
    \item choose an optimization algorithm;
    \item select line-search or trust-region globalization;
    \item verify derivative calculations;
    \item monitor objective values and constraint violations;
    \item monitor gradient or KKT residuals;
    \item inspect step sizes and trust-region radii;
    \item exploit sparsity;
    \item solve internal linear systems accurately enough;
    \item check stopping criteria;
    \item test sensitivity to initialization;
    \item verify second-order behavior when relevant;
    \item compare with alternative methods or initial guesses;
    \item independently evaluate the final solution.
\end{enumerate}

%%%%%%%%%%%%%%%%%%%%%%%%%%%%%%%%%%%%%%%%%%%%%%%%%%%%%%%%%%%%
\subsection{Choosing an Unconstrained Method}
\label{subsec:choosing-unconstrained-method}
%%%%%%%%%%%%%%%%%%%%%%%%%%%%%%%%%%%%%%%%%%%%%%%%%%%%%%%%%%%%

For inexpensive gradients and moderate dimension, consider:

\[
\boxed{
\text{BFGS}.
}
\]

For very large dimension, consider:

\[
\boxed{
\text{L-BFGS}
}
\]

or

\[
\boxed{
\text{nonlinear conjugate gradient}.
}
\]

For high-accuracy smooth problems with reliable Hessians, consider:

\[
\boxed{
\text{Newton or trust-region Newton methods}.
}
\]

For large data sets, consider:

\[
\boxed{
\text{stochastic first-order methods}.
}
\]

%%%%%%%%%%%%%%%%%%%%%%%%%%%%%%%%%%%%%%%%%%%%%%%%%%%%%%%%%%%%
\subsection{Choosing a Constrained Method}
\label{subsec:choosing-constrained-method}
%%%%%%%%%%%%%%%%%%%%%%%%%%%%%%%%%%%%%%%%%%%%%%%%%%%%%%%%%%%%

For simple convex sets, consider:

\[
\boxed{
\text{projected-gradient or proximal methods}.
}
\]

For smooth nonlinear constraints, consider:

\[
\boxed{
\text{SQP}
}
\]

or

\[
\boxed{
\text{interior-point methods}.
}
\]

For equality constraints or difficult feasibility, consider:

\[
\boxed{
\text{augmented Lagrangian methods}.
}
\]

For bound-constrained problems, specialized active-set or projected methods
can be highly effective.

%%%%%%%%%%%%%%%%%%%%%%%%%%%%%%%%%%%%%%%%%%%%%%%%%%%%%%%%%%%%
\subsection{Common Errors}
\label{subsec:numerical-optimization-common-errors}
%%%%%%%%%%%%%%%%%%%%%%%%%%%%%%%%%%%%%%%%%%%%%%%%%%%%%%%%%%%%

\subsubsection{Using a Search Direction That Is Not Descent}

A line search cannot guarantee decrease if

\[
\boxed{
\nabla f(x_k)^Tp_k\geq0.
}
\]

The direction itself must first be appropriate.

%%%%%%%%%%%%%%%%%%%%%%%%%%%%%%%%%%%%%%%%%%%%%%%%%%%%%%%%%%%%
\subsubsection{Using a Step Size That Is Always Too Large}

A large fixed step may cause:

\[
\boxed{
\text{oscillation},
}
\]

\[
\boxed{
\text{divergence},
}
\]

or

\[
\boxed{
\text{constraint violations}.
}
\]

Line search or adaptive step selection can improve robustness.

%%%%%%%%%%%%%%%%%%%%%%%%%%%%%%%%%%%%%%%%%%%%%%%%%%%%%%%%%%%%
\subsubsection{Using a Step Size That Is Always Too Small}

Extremely small steps may produce apparent stability but very slow
progress.

A tiny objective change does not automatically indicate convergence.

%%%%%%%%%%%%%%%%%%%%%%%%%%%%%%%%%%%%%%%%%%%%%%%%%%%%%%%%%%%%
\subsubsection{Forming a Hessian Inverse Explicitly}

Newton's direction should normally be computed by solving

\[
\boxed{
H_kp_k=-g_k.
}
\]

Explicitly forming

\[
H_k^{-1}
\]

is generally unnecessary.

%%%%%%%%%%%%%%%%%%%%%%%%%%%%%%%%%%%%%%%%%%%%%%%%%%%%%%%%%%%%
\subsubsection{Using Newton's Method With an Indefinite Hessian Without Safeguards}

The Newton direction may point uphill.

Use:

\[
\boxed{
\text{Hessian modification},
}
\]

\[
\boxed{
\text{line search},
}
\]

or

\[
\boxed{
\text{trust regions}.
}
\]

%%%%%%%%%%%%%%%%%%%%%%%%%%%%%%%%%%%%%%%%%%%%%%%%%%%%%%%%%%%%
\subsubsection{Assuming a Small Gradient Proves a Minimum}

A stationary point may be:

\[
\boxed{
\text{minimum},
}
\]

\[
\boxed{
\text{maximum},
}
\]

or

\[
\boxed{
\text{saddle point}.
}
\]

Second-order information may be needed.

%%%%%%%%%%%%%%%%%%%%%%%%%%%%%%%%%%%%%%%%%%%%%%%%%%%%%%%%%%%%
\subsubsection{Ignoring Conditioning}

A poorly conditioned objective can make gradient descent extremely slow.

Scaling or curvature-aware methods may be necessary.

%%%%%%%%%%%%%%%%%%%%%%%%%%%%%%%%%%%%%%%%%%%%%%%%%%%%%%%%%%%%
\subsubsection{Using Finite Differences With Arbitrary Step Size}

Derivative quality depends critically on \(h\).

Too large creates truncation error.

Too small creates cancellation and roundoff error.

%%%%%%%%%%%%%%%%%%%%%%%%%%%%%%%%%%%%%%%%%%%%%%%%%%%%%%%%%%%%
\subsubsection{Trusting Unverified Gradients}

A sign or indexing error in a gradient can completely invalidate
optimization results.

Gradient checking should be performed before blaming the solver.

%%%%%%%%%%%%%%%%%%%%%%%%%%%%%%%%%%%%%%%%%%%%%%%%%%%%%%%%%%%%
\subsubsection{Assuming Automatic Differentiation Fixes the Model}

AD computes derivatives of the implemented computation.

If the model or code is wrong, the derivative is accurately wrong.

%%%%%%%%%%%%%%%%%%%%%%%%%%%%%%%%%%%%%%%%%%%%%%%%%%%%%%%%%%%%
\subsubsection{Ignoring Constraint Violation}

A low objective value is meaningless if the final point is infeasible.

Always report constraint residuals.

%%%%%%%%%%%%%%%%%%%%%%%%%%%%%%%%%%%%%%%%%%%%%%%%%%%%%%%%%%%%
\subsubsection{Using Huge Penalty Parameters Immediately}

Very large penalty coefficients can create severe numerical ill
conditioning.

Increase penalties systematically or use augmented Lagrangian methods.

%%%%%%%%%%%%%%%%%%%%%%%%%%%%%%%%%%%%%%%%%%%%%%%%%%%%%%%%%%%%
\subsubsection{Ignoring KKT Residuals}

For constrained optimization, a small gradient of the objective alone does
not indicate optimality.

Stationarity must include constraint multipliers.

%%%%%%%%%%%%%%%%%%%%%%%%%%%%%%%%%%%%%%%%%%%%%%%%%%%%%%%%%%%%
\subsubsection{Using One Stopping Criterion Alone}

A solver may satisfy a small-step test because it has stalled.

Use multiple measures such as:

\[
\boxed{
\text{gradient},
}
\]

\[
\boxed{
\text{step},
}
\]

\[
\boxed{
\text{objective change},
}
\]

and

\[
\boxed{
\text{feasibility}.
}
\]

%%%%%%%%%%%%%%%%%%%%%%%%%%%%%%%%%%%%%%%%%%%%%%%%%%%%%%%%%%%%
\subsubsection{Ignoring Scale in Tolerances}

Absolute tolerances can be inappropriate when variables or objectives have
very large or very small magnitudes.

Relative tolerances are often more meaningful.

%%%%%%%%%%%%%%%%%%%%%%%%%%%%%%%%%%%%%%%%%%%%%%%%%%%%%%%%%%%%
\subsubsection{Assuming SGD Should Decrease the Objective Every Step}

Stochastic gradients contain sampling noise.

Individual steps may increase the full objective even while the method
makes progress in expectation.

%%%%%%%%%%%%%%%%%%%%%%%%%%%%%%%%%%%%%%%%%%%%%%%%%%%%%%%%%%%%
\subsubsection{Assuming One Initial Point Is Enough for Nonconvex Problems}

A local algorithm may converge to different stationary points from
different starts.

Multiple starts can reveal sensitivity to initialization.

%%%%%%%%%%%%%%%%%%%%%%%%%%%%%%%%%%%%%%%%%%%%%%%%%%%%%%%%%%%%
\subsubsection{Ignoring the Cost of Linear Algebra Inside Optimization}

For Newton, SQP, and interior-point methods, solving linear systems may
dominate runtime.

Optimization efficiency therefore depends heavily on numerical linear
algebra.

%%%%%%%%%%%%%%%%%%%%%%%%%%%%%%%%%%%%%%%%%%%%%%%%%%%%%%%%%%%%
\subsection{Applications}
\label{subsec:numerical-optimization-applications}
%%%%%%%%%%%%%%%%%%%%%%%%%%%%%%%%%%%%%%%%%%%%%%%%%%%%%%%%%%%%

\begin{applicationbox}

\textbf{Machine Learning.}
Training many statistical and neural-network models is a large numerical
optimization problem.

\medskip

\textbf{Statistics.}
Maximum likelihood, nonlinear regression, penalized estimation, and
Bayesian approximations rely on numerical optimization.

\medskip

\textbf{Engineering Design.}
Geometry, material use, performance, energy, and safety constraints are
optimized numerically.

\medskip

\textbf{Optimal Control.}
Direct transcription converts trajectory optimization into large nonlinear
programs.

\medskip

\textbf{Inverse Problems.}
Unknown model parameters are estimated by minimizing mismatch between
predictions and measurements.

\medskip

\textbf{Finance.}
Portfolio allocation, calibration, risk minimization, and parameter
estimation require constrained numerical optimization.

\medskip

\textbf{Scientific Computing.}
Many PDE, simulation, and parameter-identification problems produce
large-scale nonlinear optimization models.

\medskip

\textbf{Operations Research.}
Continuous relaxations, nonlinear programs, and large sparse constrained
models depend on robust numerical solvers.

\end{applicationbox}

%%%%%%%%%%%%%%%%%%%%%%%%%%%%%%%%%%%%%%%%%%%%%%%%%%%%%%%%%%%%
\subsection{Exercises}
\label{subsec:numerical-optimization-exercises}
%%%%%%%%%%%%%%%%%%%%%%%%%%%%%%%%%%%%%%%%%%%%%%%%%%%%%%%%%%%%

\begin{exercise}[1]
Define a descent direction.
\end{exercise}

\begin{exercise}[1]
Define a line search.
\end{exercise}

\begin{exercise}[1]
Define gradient descent.
\end{exercise}

\begin{exercise}[1]
Define a Newton step.
\end{exercise}

\begin{exercise}[1]
Define a quasi-Newton method.
\end{exercise}

\begin{exercise}[1]
Define a trust region.
\end{exercise}

\begin{exercise}[1]
Define a penalty method.
\end{exercise}

\begin{exercise}[1]
Define a barrier method.
\end{exercise}

\begin{exercise}[1]
Define an augmented Lagrangian.
\end{exercise}

\begin{exercise}[1]
Define stochastic gradient descent.
\end{exercise}

\begin{exercise}[2]
Determine whether a given vector \(p\) is a descent direction by evaluating

\[
\nabla f(x)^Tp.
\]
\end{exercise}

\begin{exercise}[2]
Perform one gradient-descent step for a quadratic objective.
\end{exercise}

\begin{exercise}[2]
Test whether a given trial step satisfies the Armijo condition.
\end{exercise}

\begin{exercise}[3]
Perform a complete backtracking line search for a specified function and
direction.
\end{exercise}

\begin{exercise}[3]
Check the Wolfe conditions for a given step.
\end{exercise}

\begin{exercise}[3]
Derive the exact steepest-descent step for a positive-definite quadratic.
\end{exercise}

\begin{exercise}[3]
Explain why ill conditioning causes zigzagging in steepest descent.
\end{exercise}

\begin{exercise}[2]
Construct the gradient and Hessian of a given objective.
\end{exercise}

\begin{exercise}[3]
Compute one Newton step.
\end{exercise}

\begin{exercise}[3]
Determine whether the Hessian is positive definite.
\end{exercise}

\begin{exercise}[3]
Determine whether the Newton direction is a descent direction.
\end{exercise}

\begin{exercise}[3]
Modify an indefinite Hessian using

\[
H+\tau I.
\]
\end{exercise}

\begin{exercise}[3]
Apply damped Newton with a specified line-search step.
\end{exercise}

\begin{exercise}[3]
Compute

\[
s_k=x_{k+1}-x_k
\]

and

\[
y_k=g_{k+1}-g_k.
\]
\end{exercise}

\begin{exercise}[3]
Verify the secant condition for a given quasi-Newton approximation.
\end{exercise}

\begin{exercise}[3]
Perform one BFGS update on a small problem.
\end{exercise}

\begin{exercise}[3]
Verify symmetry of the BFGS update.
\end{exercise}

\begin{exercise}[3]
Explain why

\[
y_k^Ts_k>0
\]

is important for BFGS.
\end{exercise}

\begin{exercise}[3]
Explain why L-BFGS is useful in high dimensions.
\end{exercise}

\begin{exercise}[3]
Compute a Fletcher--Reeves nonlinear conjugate-gradient coefficient.
\end{exercise}

\begin{exercise}[3]
Compute a Polak--Ribiere coefficient.
\end{exercise}

\begin{exercise}[2]
Write the quadratic trust-region subproblem.
\end{exercise}

\begin{exercise}[3]
Compute actual and predicted reductions.
\end{exercise}

\begin{exercise}[3]
Compute the trust-region ratio

\[
\rho_k.
\]
\end{exercise}

\begin{exercise}[3]
Decide whether the step should be accepted and whether the radius should
increase or decrease.
\end{exercise}

\begin{exercise}[3]
Find the Cauchy point for a simple trust-region model.
\end{exercise}

\begin{exercise}[2]
Approximate a gradient using forward finite differences.
\end{exercise}

\begin{exercise}[3]
Approximate the same gradient using centered differences.
\end{exercise}

\begin{exercise}[3]
Compare finite-difference gradients with an exact gradient.
\end{exercise}

\begin{exercise}[3]
Explain forward-mode automatic differentiation.
\end{exercise}

\begin{exercise}[3]
Explain reverse-mode automatic differentiation.
\end{exercise}

\begin{exercise}[3]
Explain why reverse mode is attractive for scalar objectives with many
inputs.
\end{exercise}

\begin{exercise}[3]
Perform a directional derivative gradient check.
\end{exercise}

\begin{exercise}[2]
Project a point onto a box constraint.
\end{exercise}

\begin{exercise}[3]
Perform one projected-gradient step.
\end{exercise}

\begin{exercise}[3]
Construct a quadratic penalty function for an equality constraint.
\end{exercise}

\begin{exercise}[3]
Construct a logarithmic barrier for two inequality constraints.
\end{exercise}

\begin{exercise}[3]
Construct an augmented Lagrangian.
\end{exercise}

\begin{exercise}[3]
Perform one multiplier update.
\end{exercise}

\begin{exercise}[3]
Write an SQP subproblem for a simple nonlinear constrained problem.
\end{exercise}

\begin{exercise}[3]
Construct a merit function.
\end{exercise}

\begin{exercise}[3]
Explain the difference between active-set and interior-point methods.
\end{exercise}

\begin{exercise}[3]
Construct stationarity, feasibility, and complementarity residuals.
\end{exercise}

\begin{exercise}[3]
Design a stopping test for an unconstrained optimizer.
\end{exercise}

\begin{exercise}[3]
Design a stopping test for a constrained optimizer.
\end{exercise}

\begin{exercise}[3]
Scale a problem whose variables differ by six orders of magnitude.
\end{exercise}

\begin{exercise}[3]
Explain preconditioning of gradient descent.
\end{exercise}

\begin{exercise}[3]
Write a Newton--KKT system for an equality-constrained problem.
\end{exercise}

\begin{exercise}[3]
Derive a Schur complement for a block KKT system.
\end{exercise}

\begin{exercise}[3]
Explain how Hessian-vector products enable matrix-free optimization.
\end{exercise}

\begin{exercise}[3]
Write an inexact Newton residual condition.
\end{exercise}

\begin{exercise}[2]
Compute a subgradient of

\[
f(x)=|x|
\]

at a nonzero point.
\end{exercise}

\begin{exercise}[3]
Describe the subdifferential of

\[
|x|
\]

at \(x=0\).
\end{exercise}

\begin{exercise}[3]
Perform one subgradient step.
\end{exercise}

\begin{exercise}[3]
Compute the proximal operator of a simple quadratic function.
\end{exercise}

\begin{exercise}[3]
Apply soft thresholding to a vector.
\end{exercise}

\begin{exercise}[3]
Perform one proximal-gradient step for an \(\ell_1\)-regularized problem.
\end{exercise}

\begin{exercise}[2]
Write the full gradient of a finite-sum objective.
\end{exercise}

\begin{exercise}[2]
Write one SGD update.
\end{exercise}

\begin{exercise}[3]
Construct a mini-batch gradient.
\end{exercise}

\begin{exercise}[3]
Explain why an SGD step may increase the full objective.
\end{exercise}

\begin{exercise}[3]
Check whether a sequence of step sizes satisfies

\[
\sum\alpha_k=\infty
\]

and

\[
\sum\alpha_k^2<\infty.
\]
\end{exercise}

\begin{exercise}[3]
Perform one momentum update.
\end{exercise}

\begin{exercise}[3]
Explain the difference between deterministic and stochastic optimization
iterations.
\end{exercise}

\begin{exercise}[3]
Explain why multiple initializations may be necessary for a nonconvex
objective.
\end{exercise}

\begin{exercise}[4]
Prove descent of the negative-gradient direction.
\end{exercise}

\begin{exercise}[4]
Study global convergence of gradient descent under Lipschitz-gradient
assumptions.
\end{exercise}

\begin{exercise}[4]
Derive convergence rates for gradient descent on strongly convex
functions.
\end{exercise}

\begin{exercise}[4]
Derive steepest-descent convergence for positive-definite quadratics.
\end{exercise}

\begin{exercise}[4]
Prove local quadratic convergence of Newton's method for optimization.
\end{exercise}

\begin{exercise}[4]
Study damped Newton globalization.
\end{exercise}

\begin{exercise}[4]
Derive the BFGS update from secant and minimum-change principles.
\end{exercise}

\begin{exercise}[4]
Study superlinear convergence of BFGS.
\end{exercise}

\begin{exercise}[4]
Study limited-memory quasi-Newton methods.
\end{exercise}

\begin{exercise}[4]
Study nonlinear conjugate-gradient global convergence.
\end{exercise}

\begin{exercise}[4]
Derive the Cauchy-decrease bound for trust-region methods.
\end{exercise}

\begin{exercise}[4]
Study the trust-region optimality conditions.
\end{exercise}

\begin{exercise}[4]
Study dogleg methods.
\end{exercise}

\begin{exercise}[4]
Study truncated conjugate gradient for trust-region subproblems.
\end{exercise}

\begin{exercise}[4]
Study convergence of projected-gradient methods.
\end{exercise}

\begin{exercise}[4]
Study exact penalty functions.
\end{exercise}

\begin{exercise}[4]
Derive first-order conditions for logarithmic barrier problems.
\end{exercise}

\begin{exercise}[4]
Study augmented Lagrangian convergence.
\end{exercise}

\begin{exercise}[4]
Study local SQP convergence.
\end{exercise}

\begin{exercise}[4]
Derive primal-dual Newton equations for interior-point methods.
\end{exercise}

\begin{exercise}[4]
Study filter globalization for constrained optimization.
\end{exercise}

\begin{exercise}[4]
Study semismooth Newton methods.
\end{exercise}

\begin{exercise}[4]
Study convergence of proximal-gradient methods.
\end{exercise}

\begin{exercise}[4]
Study accelerated proximal-gradient methods.
\end{exercise}

\begin{exercise}[4]
Study stochastic-gradient convergence under unbiased-noise assumptions.
\end{exercise}

\begin{exercise}[4]
Study variance reduction for finite-sum optimization.
\end{exercise}

\begin{exercise}[4]
Investigate optimization with automatic differentiation and
Hessian-vector products.
\end{exercise}

%%%%%%%%%%%%%%%%%%%%%%%%%%%%%%%%%%%%%%%%%%%%%%%%%%%%%%%%%%%%
\subsection{Conceptual Summary}
\label{subsec:numerical-optimization-summary}
%%%%%%%%%%%%%%%%%%%%%%%%%%%%%%%%%%%%%%%%%%%%%%%%%%%%%%%%%%%%

Numerical optimization constructs iterates

\[
\boxed{
x_{k+1}
=
x_k+\alpha_kp_k
}
\]

that approach an optimal or stationary point.

Gradient descent uses

\[
\boxed{
p_k
=
-\nabla f(x_k).
}
\]

Newton's method solves

\[
\boxed{
\nabla^2f(x_k)p_k
=
-\nabla f(x_k).
}
\]

Quasi-Newton methods approximate curvature using

\[
\boxed{
s_k
=
x_{k+1}-x_k
}
\]

and

\[
\boxed{
y_k
=
g_{k+1}-g_k.
}
\]

Trust-region methods solve

\[
\boxed{
\min_{\|p\|\leq\Delta_k}
m_k(p).
}
\]

Constrained numerical optimization uses methods such as

\[
\boxed{
\text{projection},
}
\]

\[
\boxed{
\text{penalty},
}
\]

\[
\boxed{
\text{barrier},
}
\]

\[
\boxed{
\text{augmented Lagrangian},
}
\]

\[
\boxed{
\text{SQP},
}
\]

and

\[
\boxed{
\text{interior point}.
}
\]

Large-scale problems rely on:

\[
\boxed{
\text{sparsity},
}
\]

\[
\boxed{
\text{Hessian-vector products},
}
\]

\[
\boxed{
\text{iterative linear algebra},
}
\]

and sometimes

\[
\boxed{
\text{stochastic gradients}.
}
\]

Reliable numerical optimization therefore combines

\[
\boxed{
\text{optimization theory}
+
\text{numerical analysis}
+
\text{linear algebra}
+
\text{algorithm design}
+
\text{computational verification}.
}
\]

%%%%%%%%%%%%%%%%%%%%%%%%%%%%%%%%%%%%%%%%%%%%%%%%%%%%%%%%%%%%
\subsection{Section Connections}
\label{subsec:numerical-optimization-connections}
%%%%%%%%%%%%%%%%%%%%%%%%%%%%%%%%%%%%%%%%%%%%%%%%%%%%%%%%%%%%

\chapterconnection{
Numerical Optimization combines the mathematical optimality theory of
Chapter~10 with the conditioning, stability, and linear-algebra tools of
Sections~\ref{sec:numerical-analysis} and
\ref{sec:numerical-linear-algebra}. Gradient, Newton, quasi-Newton,
trust-region, SQP, interior-point, proximal, and stochastic methods provide
computational mechanisms for solving problems that rarely admit closed
forms. These methods also appear inside optimal-control solvers, inverse
problems, statistical estimation, machine learning, and scientific
simulation. The next section, Numerical ODEs, develops computational
methods for approximating solutions of ordinary differential equations,
including Euler methods, Runge--Kutta schemes, multistep methods,
stability, stiffness, and adaptive time stepping.
}

%%%%%%%%%%%%%%%%%%%%%%%%%%%%%%%%%%%%%%%%%%%%%%%%%%%%%%%%%%%%
% SECTION SOURCE TRACKING
%%%%%%%%%%%%%%%%%%%%%%%%%%%%%%%%%%%%%%%%%%%%%%%%%%%%%%%%%%%%

%------------------------------------------------------------
% 11.3 Numerical Optimization
%------------------------------------------------------------
%
% Source basis:
%   - Standard numerical nonlinear optimization.
%   - Standard unconstrained and constrained optimization algorithms.
%   - Standard large-scale, nonsmooth, and stochastic optimization
%     techniques.
%
% Used for:
%   - iterative optimization;
%   - descent directions;
%   - gradient descent;
%   - line search;
%   - exact and inexact line search;
%   - Armijo condition;
%   - Wolfe and strong Wolfe conditions;
%   - quadratic objectives;
%   - conditioning and steepest descent;
%   - Newton's method;
%   - damped and modified Newton methods;
%   - quasi-Newton methods;
%   - secant conditions;
%   - BFGS and inverse BFGS;
%   - DFP preview;
%   - L-BFGS;
%   - nonlinear conjugate gradient;
%   - Fletcher--Reeves and Polak--Ribiere formulas;
%   - trust-region methods;
%   - trust-region ratio and radius updates;
%   - Cauchy point;
%   - derivative-free optimization;
%   - coordinate and simplex-type search;
%   - finite-difference gradients;
%   - automatic differentiation;
%   - gradient checking;
%   - projected-gradient methods;
%   - penalty methods;
%   - barrier methods;
%   - augmented Lagrangians;
%   - SQP;
%   - merit and filter methods;
%   - active-set methods;
%   - interior-point methods;
%   - primal-dual complementarity;
%   - KKT residuals;
%   - stopping criteria;
%   - scaling and preconditioning;
%   - sparse optimization;
%   - Newton--KKT systems;
%   - Schur complements;
%   - Hessian-vector products;
%   - Newton--CG and inexact Newton;
%   - nonsmooth and subgradient methods;
%   - proximal operators;
%   - proximal gradient;
%   - soft thresholding;
%   - stochastic-gradient methods;
%   - mini-batches;
%   - stochastic step sizes;
%   - momentum;
%   - acceleration and adaptive methods preview;
%   - local and global optimization;
%   - saddle points and negative curvature;
%   - warm starts and continuation;
%   - applications and exercises.
%
% Notes:
%   - Numerical ODEs are developed next in Section 11.4.
%   - Numerical PDEs follow in Section 11.5.
%   - Monte Carlo Methods follow in Section 11.6.
%   - Numerical Linear Algebra appears in Section 11.2.
%
%%%%%%%%%%%%%%%%%%%%%%%%%%%%%%%%%%%%%%%%%%%%%%%%%%%%%%%%%%%%